
\documentclass[journal=esthag,manuscript=article]{achemso}


\usepackage[version=3]{mhchem} % Formula subscripts using \ce{}
\usepackage{subcaption}
\usepackage{booktabs}          % Better tables
\usepackage{graphicx}          % Images
\usepackage{amsmath}           % Math
\usepackage{amssymb}           % \mathbb, used in the Supporting Information
\usepackage{placeins}
\renewcommand{\topfraction}{0.9}
\renewcommand{\bottomfraction}{0.9}
\renewcommand{\textfraction}{0.1}        % allow a page to be mostly float
\renewcommand{\floatpagefraction}{0.8}
\setcounter{topnumber}{3}
\setcounter{totalnumber}{5}


\author{Tadeusz Price}
\affiliation[RIVM]{National Institute for Public Health and the Environment, Bilthoven, The Netherlands.}


\author{Oliver Urs Lenz}
\affiliation[SDM]{Department of Statistics, Data Science and Modelling, National Institute for Public Health and the Environment, Bilthoven, The Netherlands.}
\alsoaffiliation[LIACS]{Leiden Institute of Advanced Computer Science, Leiden University, Leiden, the Netherlands}
\alsoaffiliation[WINST]{Department of Mathematics, Computer Science and Statistics, Ghent University, Ghent, Belgium}

\author{Leo Posthuma}
\affiliation[DMG]{Centre for Sustainability, Environment and Health, National Institute for Public Health and the Environment, Bilthoven, The Netherlands.}
\alsoaffiliation[Nijmegen]{Department of Environmental Science, Radboud University, Nijmegen, The Netherlands}

\title{Uncertainty Quantification for Aquatic Ecotoxicity Prediction}


\begin{document}

%%%%%%%%%%%%%%%%%%%%%%%%%%%%%%%%%%%%%%%%%%%%%%%%%%%%%%%%%%%%%%%%%%%%%
%% TODO Graphical abstract!!!! required by EST
%%%%%%%%%%%%%%%%%%%%%%%%%%%%%%%%%%%%%%%%%%%%%%%%%%%%%%%%%%%%%%%%%%%%%


\begin{abstract}
Aquatic ecosystems are threatened by thousands of chemicals, of which the hazard is often largely unknown. These hazard characteristics are derived from laboratory toxicity tests of individual chemicals on individual species. Part of those data consist of median Lethal Concentrations for 50\% of the test individuals (LC50), but for the overwhelming majority of chemicals in use, such data exist for only a small fraction of the species present in ecosystems. A growing body of research applies machine learning to bridge these gaps, predicting the hazard of data-poor and even untested chemicals. These approaches are promising, but before they can be relied on to set protective regulatory standards or support other decisions, apart from being accurate in general they must also quantify how far individual predictions can be trusted.

We address this by using a Bayesian modelling framework, training a Bayesian Factorization Machine (BFM) on ecotox data, which predicts not just LC50 values, but also their expected variance. This uncertainty is separated in the model into a reducible part arising from limited data (epistemic) and an irreducible part arising from inter-test and biological variability (aleatoric). The BFM achieves state-of-the-art accuracy in filling the testing gaps, but more importantly we show how its uncertainty estimates propagate into downstream risk assessment. They carry through to Species Sensitivity Distributions (SSDs) and the hazardous concentrations derived from them, indicating when predicted hazard data can and cannot be relied on for practical decision support. The same approach is not specific to ecotoxicology, and applies wherever decisions rest on sparse and noisy data.
\end{abstract}




\section{Introduction}

Aquatic ecosystems face pressure from chemical pollution. There are hundreds of thousands of chemicals in use in society \cite{wang2020toward}, and many of them find their way into surface water systems. There are increasing regulatory efforts to prevent and control this, ranging from the recently proposed \textit{Safe \& Sustainable by Design} \cite{caldeira2022safe} framework as a pro-active approach in the design of novel chemicals, up to the retrospective water quality assessment and mitigation of the EU Water Framework Directive (2000/60/EC). Reaching the societal goal of \textit{sufficient water of good quality}, for the many human uses of clean water and as habitat for aquatic species assemblages, requires the risk assessment and management of individual chemicals and chemical mixtures.

One of the key requirements in risk assessment and management is to know the hazards of the chemicals that we use. That knowledge is commonly collected by performing single-species, single-compound ecotoxicity tests in a laboratory setting. For a given chemical, standardised tests are typically performed on a handful of selected species, such as an algae, a daphnid and/or a fish. The purpose of these tests is to determine the concentration at which exposure is lethal (LCx), or has a specific negative effect on a life history characteristic (ECx), such as growth or reproduction, for x\% of the test individuals.

The results and conditions of these tests are being collated and curated by several organisations. For example, the US EPA ECOTOX Knowledgebase now contains over a million test records spanning thousands of species and chemicals \cite{ECOTOX}. However, this still represents a mere fraction all species--chemical pairs.

Due to this data sparsity, we typically only know at which concentration level a chemical is toxic for a minority of species. The traditional approach to deal with this is to fit a parametric (typically lognormal) distribution on the available single-species toxicity values of a given chemical, the so-called Species Sensitivity Distribution (SSD). This describes the sensitivity of an entire ecological community to a chemical \cite{posthuma2001species}, and can be used to derive regulatory thresholds such as the Hazardous Concentration for 5\% of species (HC5). However, this fit is only reliable when the number and variety of tested species is sufficient, which is only the case for a small fraction of all chemicals.

The growing data availability, together with the sparsity problem, has motivated a wave of machine learning approaches that use the collated testing data on known species--chemical pairs to predict the toxicity of untested chemical--species pairs, or even of entirely new chemicals or never-tested species. \citet{von2023potential} identify key parameters where machine learning could address data gaps in toxicity characterisation\iffalse OL: Can we be more specific here? What were those key parameters?\fi, \citet{kosnik2024harnessing}, advocate for computational methods as essential to characterising chemical impacts on biodiversity at scale, while \citet{dong2025machine} review the expanding role of machine learning across exposure estimation, biotoxicity prediction, and environmental fate modelling.

A more critical account is provided by \citet{schur2025comparability}, who highlight significant comparability issues in predictive ecotoxicology, calling for standardised benchmarks and evaluation protocols. To address this, the authors have curated the ADORE collection of benchmark datasets, based on the US EPA ECOTOX Knowledgebase and enriched with chemical and taxonomical features \cite{ADORE}.

The ADORE collection has enabled the systematic evaluation of various models. \citet{gasser2024fish} found that model performance was highly dependent on the data split while benchmarking tree-based and Gaussian-process regression models on the fish-specific ADORE challenge, and applying them to fish SSD construction. A novel approach was taken by \citet{Viljanen2023}, who framed ecotoxicity prediction as a recommender-system problem and showed that factorisation machines could match more complex methods in the context of highly sparse test data. This was further elaborated by \citet{PosthumaImprovingHazard}, who scaled the approach to predict LC50 values for over four million species--chemical data gaps and used the resulting predictions to construct Species Sensitivity Distributions and Chemical Hazard Distributions for all chemicals and species, representing the current state of the art for machine-learning-based SSD construction at scale. Others have explored meta-learning to share knowledge across low-resource species \cite{schlender2023bigger}, and Bio-QSARs that incorporate physiological traits to enable genuine cross-species extrapolation \cite{zubrod2024bio}.

While the overall predictive accuracy of such models can be measured using approaches like cross-validation, yielding a single accuracy metric across the entire dataset, this tells us little about how much to trust an \textit{individual} prediction (and by extension resulting scatter-SSDs and derived protective thresholds). This is because the  existing models produce only point estimates, and do not quantify the confidence of individual predictions. This is a significant limitation for practical risk assessment, where knowing \textit{how much} to trust a prediction is as important as the prediction itself. In practice, we expect prediction quality to be impacted by two sources of uncertainty. Firstly, a chemical that has been tested on dozens of species should yield more reliable predictions than one with only a handful of observations, and similarly for well-studied versus rarely tested species. Secondly, the currently available data clearly show, for species--chemical pairs that have been tested multiple times, that repeated testing of the chemical on the same species yields highly variable results. These two sources of uncertainty are commonly referred to as epistemic (arising from limited data constraining the model) and aleatoric (inherent to the biological and experimental variability of the testing process) \cite{epistemicaleatoric}.

Bayesian modelling offers a natural framework to deal with uncertainty: by treating model parameters as distributions rather than fixed values, it produces not just a single prediction but a range of plausible outcomes, directly expressing the model's uncertainty. Bayesian approaches have recently been applied to uncertainty-aware human toxicity prediction \cite{von2026uncertainty} and to hierarchical SSD construction for microplastic particles \cite{iwasaki2026bayesian}, but have not yet been applied to the large-scale species--chemical matrix completion problem in ecotoxicology.

The present work addresses this gap by extending the factorisation machine architecture applied to ecotoxicity prediciton in \citet{PosthumaImprovingHazard} into a fully Bayesian framework, the Bayesian Factorisation Machine (BFM), adapted from \citet{rendle2011bayesian}. In addition, we explicitly model how noisy the underlying data is for each chemical, allowing us to distinguish between aleatoric and epistemic uncertainty.

Thus, the BFM model presented in this work allows risk assessors to see predicted toxicity values for all combinations of species and chemicals that are present in the training data, the uncertainty surrounding these predictions, as well as to how much of this uncertainty is inherent experimental noise, and how much is due to a lack of data.

Furthermore, we also describe how this uncertainty propagates to the next steps of practical assessments. For each chemical we derive an SSD that is made up of the point-predictions for each species, along with an uncertainty band showing the space of plausible SSD curves given the model uncertainty, allowing for more informed risk assessment and regulation.

In this study, the LC50 endpoint is used as a methodological prototype: acute lethality is a relatively well-defined test endpoint (compared to e.g.\ chronic no-observed effect concentrations for life-history characteristics such as growth and reproduction) and is numerically abundant in the data, making it ideal for developing and validating the uncertainty quantification approach. However, with additional tuning and validation, the method can easily be adapted to other test endpoints such as NOEC or EC10 under chronic exposure conditions, which are often preferred as basis for regulatory applications (such as setting protective concentration standards).


\section{Materials and Methods}
\subsection{Overview}

This section sets out the data and the model. We begin with the ecotoxicity dataset and the three properties that shape the modelling problem: the data are \textbf{sparse}, they are \textbf{unevenly distributed} across chemicals and species, and they show substantial \textbf{variability} between repeated tests of the same chemical and species. Each property motivates a specific modelling choice, which we make explicit as we go.

The prediction task is to estimate the LC50 for species--chemical combinations that have not been tested, using the combinations that have. Our starting point is the Factorization Machine applied to this task by \citet{PosthumaImprovingHazard}, which is accurate and well suited to sparse data. We keep its predictive core unchanged and make two changes that together turn a single point prediction into a prediction with quantified, decomposed uncertainty.

The first change is to reformulate the model in Bayesian terms. Instead of fitting one best set of parameters, we treat the parameters as quantities that the data constrain only partially, and we work with many plausible parameter sets at once. Each set produces its own prediction, so every species--chemical combination receives not a single number but a spread of predictions. The width of that spread is a per-prediction measure of how strongly the data constrain the model, which is information a point estimate cannot provide. Intuitively, the data are consistent with many models, and we account for that rather than committing to one.

The second change is to model the test noise explicitly, separately for each chemical. The repeated-test variability in the data shows that a single laboratory result is better understood as one draw from a distribution than as a fixed value. We therefore add a variance term that the model learns for each chemical, on the assumption that different chemicals carry different baseline levels of test noise. This lets the model report how much of the spread around a prediction is irreducible measurement noise rather than a lack of data.

The remainder of the section describes the dataset, then the predictive model and these two changes in turn. The full mathematical specification, including the hierarchical structure and the sampling procedure, is given in the Supporting Information. All code is available online, allowing the results of this paper to be reproduced, and the method to be applied to other datasets and adapted as desired. 


\subsection{Dataset}

The dataset used is an unfiltered view of the ADORE dataset introduced by \citet{ADORE}. The ecotoxicity values in ADORE come from the US EPA ECOTOX dataset \citep{ECOTOX}, and ADORE includes a large number of features of chemicals and species, forming a low barrier to entry for data scientists from outside the domain and reducing the risk of errors when procuring features. Multiple endpoints are available, such as Effective Concentration or No Effect Concentration, however out of the available endpoints, LC50 is chosen as the sole endpoint to focus on for all experiments, as it is consistently well-defined and has a significant amount of data available. Future work could use data from multiple endpoints to improve model performance, however this would require significant data engineering work to standardise the reporting of other endpoints.

Each row of the dataset contains an LC50 observation for a tested species--chemical pair in weight-based concentration units (mg/L), which is log-transformed for modelling. Each observation includes information relating to the chemical that was tested, the species that was exposed, and further information about the lab conditions (such as test duration, and various other test characteristics, if available).

The ADORE dataset includes a number of `challenge' slices, where the data is limited to specific sub-species, however, in order to leverage as much of the data as possible, the processed but not divided dataset is used here. This is the same unfiltered view of ADORE used in \citet{PosthumaImprovingHazard}, so that the exact same results as presented in that paper are expanded with the uncertainty analyses, allowing for direct comparisons. Specifically \texttt{ecotox\_mortality\_processed.csv} is joined with \texttt{ecotox\_properties\_with-oecd-function.csv}, available to download at \url{https://gitlab.renkulab.io/mltox/adore}, where in-depth details on the dataset processing can also be found. Table~\ref{tab:adore_summary} shows summary statistics about the dataset.

\begin{table}[ht]
\centering
\caption{Summary statistics of the ADORE (mortality processed) dataset}
\label{tab:adore_summary}
\begin{tabular}{l r}
\hline
\textbf{Statistic} & \textbf{Value} \\
\hline
Number of unique chemicals & 3,295 \\
Number of unique species & 1,267 \\
Number of durations & 4 \\
Number of unique (species, chemical, duration) triplets observed & 31,319 \\
Total number of observations & 70,670 \\
Sparsity (species$\times$chemical matrix) & 99.5\% \\
\hline
\end{tabular}
\end{table}

The overview of ADORE (ecotox\_mortality\_processed) in Table~\ref{tab:adore_summary} illustrates that ecotoxicity data sets have specific characteristics, posing challenges for practice. The Table shows the first major challenge of this dataset: \textbf{sparsity}. Even if the duration dimension (with four exposure times: 24, 48, 72 and 96 hours of exposure) is collapsed, only around 0.5\% of the chemical--species pairs have observed LC50 values. This sparsity is the reason a predictive model is needed at all: for the large majority of pairs there is no measured value, so the task is to infer those values from the few pairs that have been tested. Note further that the data set is not only sparse, but also just a fraction of the full matrix, which would need to cover more than 350,000 chemicals that are in active use globally \cite{wang2020toward}, and a vastly larger diversity of species than currently in ADORE.

The second major challenge for using the data effectively is the \textbf{uneven distribution of observations} across the matrix, which is illustrated in Figure~\ref{fig:rank-freq} from the points of view of chemicals and species separately. The rank-frequency plots illustrate that most unique species and chemicals have just a handful of observations, while a few chemicals and species have been tested hundreds of times. As a result, the amount of evidence behind a prediction varies widely across the matrix: a prediction involving a chemical or species with hundreds of observations is far better supported than one involving a chemical or species seen only a few times. We will want the model's reported confidence to track this difference.

\begin{figure}[htbp]
    \centering
    \includegraphics[width=1\linewidth]{images/rank_freq.png}
    \caption{Characteristics of the species--chemical pair test data in the ADORE data set: plots showing the number of observations per species (left) and per chemical (right) ranked by number of observations.}
    \label{fig:rank-freq}
\end{figure}

The third challenge of the data is very high \textbf{inter-test variability} among the multiple-tested species--chemical pairs. Figure~\ref{fig:rsd-distr} shows the distribution of normalised standard deviations for chemical--species--duration combinations where at least 10 observations are available. This shows that there is a significant portion of the data set where the variance of repeated observations is very high, greater than the mean, or even double or triple the mean. Some of this spread could in principle be reduced: more standardised test conditions and protocols would remove part of it, and, on the modelling side, additional test-condition features could explain part of it as structured effects rather than leaving it as unexplained noise. However far such efforts are taken, though, a residual component will always remain, rooted in the biological and experimental variability inherent to toxicity testing. Rather than treating this as a nuisance to be minimised, our approach confronts it directly by quantifying it, so that whatever steps the community takes to reduce inter-test variability, accounting for what remains is not sidestepped. Because the same nominal test can give very different results, a single measured LC50 is best understood as one sample from an underlying distribution rather than a fixed value. A model that returns only a point prediction cannot represent this spread, which is one reason point predictions alone are not sufficient here; we return to this when we model the noise explicitly.


These three properties: sparsity, uneven sampling, and high inter-test variability, are not unique to the ADORE data but are common to ecotoxicity datasets, and they drive model choice and evaluation in applying machine learning in this domain. Each one motivates a specific element of the model that follows: sparsity motivates a model that can fill gaps, uneven sampling motivates a per-prediction measure of confidence, and inter-test variability motivates an explicit model of the test noise.

\begin{figure}[htbp]
    \centering
    \includegraphics[width=0.7\linewidth]{images/RSD_distribution_all_durations.png}
    \caption{Characteristics of the species--chemical--duration sets of test data in the ADORE data set, for which there are more than 10 observations for each combination (multiple test data available for same triplet): the distribution of Relative Standard Deviations (RSDs) in the data grouped by species, chemical, duration illustrates the wide variability ranges of tests that are, in principle, the same.}
    \label{fig:rsd-distr}
\end{figure}
\subsection{The Hierarchical Bayesian Factorization Machine}

This section describes the model that predicts LC50 values for untested species--chemical combinations, together with the uncertainty attached to each prediction. It applies to combinations whose chemical and species each appear somewhere in the training data, and is not designed for chemicals or species that are entirely absent.

The predictive core is the factorization machine applied to this task by \citet{PosthumaImprovingHazard}, chosen because it handles sparse data well and has proven successful both on LC50 prediction and on sparse matrix completion more generally. On its own it supplies the point prediction: a single best estimate of toxicity for each combination. We keep this core unchanged and cast it in Bayesian form, following the Bayesian Factorization Machine of \citet{rendle2011bayesian}; replacing the single fitted model with an ensemble of plausible fits yields the epistemic uncertainty, the reducible part that reflects how strongly the data constrain each prediction. Working within that same framework, we implement the model ourselves and adapt it to ecotoxicity data: in place of the single global noise level assumed by the standard BFM, each chemical receives its own noise term under a shared hierarchical prior, capturing the aleatoric uncertainty, the irreducible spread between repeated tests. This per-chemical, hierarchically pooled noise model is the main methodological contribution of this work.
\begin{figure}[htbp]
    \centering
    \includegraphics[width=0.6\linewidth]{images/plate_diagram.jpg}
    \caption{Plate notation of the Hierarchical Bayesian Factorization Machine. Fixed constants (rectangles) set the hyperpriors for three groups of parameters: the linear weights $w_j$, the latent factors $v_{jf}$, and the per-chemical noise precisions $\alpha_c$. The weights $w_j$ and factors $v_{jf}$ share population-level mean and precision hyperparameters ($\mu_w, \lambda_w$ and $\mu_v, \lambda_v$ respectively), learned jointly across all $p$ features (and, for $v_{jf}$, across all $k$ latent dimensions); the noise precisions $\alpha_c$ share a rate $b_0$, learned across all $C$ chemicals, with shape $a_0$ fixed in advance. The global bias $w_0$ has a fixed prior. For each observation $i$, the predicted mean $\hat{y}_i$ (dashed, deterministic) is computed from the features $x_i$ together with $w_0$, $w_j$, and $v_{jf}$, and the observed toxicity $y_i$ is then drawn from a Gaussian centred on $\hat{y}_i$ with precision $\alpha_{c_i}$, the noise precision of the chemical $c_i$ tested in that observation. Shaded nodes are observed, open circles are latent variables inferred during fitting, and the dashed node is a deterministic function of its parents. The full equations corresponding to each node are given in the Supporting Information.}
    \label{fig:plate}
\end{figure}

Figure~\ref{fig:plate} shows the complete model in plate notation. Reading from the top, the shared hyperparameters set the prior distributions for three groups of quantities: the linear weights, the latent factors, and the per-chemical noise precisions. Each chemical draws its own noise precision from the shared noise prior, and each observed toxicity value is generated from the factorization machine mean evaluated at that observation's features, with Gaussian noise whose precision is that of the observation's chemical. The plates mark what is repeated: the features and factors that make up the parameters, the chemicals that each carry a noise precision, and the observations that make up the data. For a reader familiar with hierarchical Bayesian models the diagram specifies the model in full; the equations are collected in the Supporting Information. In the main text we keep the focus on the three architectural choices and what each one adds.

\subsubsection{The predictive model}

The quantity the model predicts is the mean log-transformed LC50 for a given species--chemical combination at a given test duration. This predictive mean is produced by a second-order factorization machine, exactly as in \citet{PosthumaImprovingHazard}:
\begin{equation}
    \hat{y}_i(\theta) = w_0 + \sum_{j=1}^{p} w_j x_{ij} + \sum_{j=1}^{p} \sum_{j' > j}^{p} \langle \mathbf{v}_j, \mathbf{v}_{j'} \rangle x_{ij} x_{ij'}
    \label{eq:fm_mean}
\end{equation}
Here $w_0$ is a global offset, the $w_j$ are weights on the individual input features, and the $\mathbf{v}_j$ are short latent vectors, one per feature, whose inner products $\langle \mathbf{v}_j, \mathbf{v}_{j'} \rangle$ give the strength of the interaction between features $j$ and $j'$. The feature vector $\mathbf{x}_i$ encodes the identity of the chemical and the species, the test duration, the taxonomic family and class of the species, and two numerical chemical properties (log-transformed molecular weight and the octanol--water partition coefficient). Full notation is given in the Supporting Information.

In effect this is a linear model with main effects and pairwise interactions, with one important difference. The interaction terms are not stored as a separate weight for every pair of features, which would be both wasteful and unusable on sparse data, but are represented in a compact low-rank form through the latent vectors. This keeps the number of parameters small and, more importantly, lets the model estimate the interaction between a chemical and a species that were never tested together, by learning a representation of each from the combinations in which they do appear. This property is what makes the factorization machine well suited to filling a sparse matrix, and is why it performs well on this task \citep{Viljanen2023, PosthumaImprovingHazard}. Because the predictions are built from the identities of the chemical and species, the model can fill untested combinations of chemicals and species that appear in the data, but it cannot predict for an entirely new chemical or species that has never been seen. 

\subsubsection{A Bayesian reformulation}

A standard factorization machine is fitted by searching for the single set of parameter values that best matches the training data, usually by minimising a loss function with gradient descent. This produces one predicted value per combination and no indication of how reliable that value is. With data that are sparse and noisy this is a real shortcoming: the model cannot be trusted equally everywhere, yet a single predicted number gives no signal of where it rests on solid evidence and where it is effectively guessing.

The Bayesian reformulation addresses this by treating the parameters not as fixed values to be optimised but as quantities that the data constrain only partially. Each parameter is given a prior distribution, which the data update into a posterior distribution describing the values that remain plausible after seeing the evidence. These priors are arranged hierarchically: related parameters, such as all the latent factors, share a common prior whose own parameters are also learned from the data. This shares information across the group, so that parameters supported by little data are pulled towards the group average rather than taking extreme values, a form of automatic regularisation known as partial pooling. We adopt the Bayesian Factorization Machine of \citet{rendle2011bayesian} for this purpose.

Training is therefore no longer a search for one best set of parameters but the characterisation of this posterior distribution. Because the posterior has no closed form, we represent it by drawing a large number of samples from it using Gibbs sampling, a Markov chain Monte Carlo method described in the Supporting Information. Gibbs sampling builds this set by updating each parameter in turn from its distribution conditional on the current values of the others. Because the chain starts from arbitrary initial values, it takes a number of iterations before these conditional updates settle into the posterior; this initial \emph{burn-in} is discarded and only the samples drawn once the chain has stabilised are retained. Each retained sample is a complete set of parameters, and so corresponds to one plausible version of the model that produces its own full set of predictions.

For any species--chemical combination we then obtain one prediction from each sample, with $\hat{y}^{(s)}(\mathbf{x}^*)$ denoting the prediction from sample $s$ computed via Equation~\ref{eq:fm_mean}. Their average is the model's point prediction,
\begin{equation}
    \hat{y}^* = \frac{1}{S} \sum_{s=1}^{S} \hat{y}^{(s)}(\mathbf{x}^*)
    \label{eq:ppm}
\end{equation}
and their spread is the \emph{epistemic} uncertainty, the part of the uncertainty that comes from the data constraining the model only partially:
\begin{equation}
    \text{Var}_{\text{epistemic}}(\mathbf{x}^*) = \frac{1}{S} \sum_{s=1}^{S} \left( \hat{y}^{(s)}(\mathbf{x}^*) - \hat{y}^* \right)^2
    \label{eq:epistemic}
\end{equation}
This spread behaves the way confidence should: where the data pin the parameters down, the sampled models agree and the spread is small; where they do not, the sampled models disagree and the spread is large. The model thus reports, for each individual prediction, how much the available evidence supports it. This is a signal a point-prediction model cannot give, and the one a risk assessor needs in order to judge how far to rely on a predicted value.

\subsubsection{Per-chemical noise}

The second change concerns the noise in the measurements themselves. The repeated-test variability documented in Figure~\ref{fig:rsd-distr} shows that testing the same chemical on the same species more than once can give substantially different results. The Bayesian reformulation above gives us the spread across plausible models given the observed data, but says nothing about this scatter between repeated tests of the same combination; it is a distinct kind of uncertainty, internal to the measurements rather than to the fit. We therefore model it directly, on the assumption that each chemical carries some inherent biological and experimental variability that should be estimated rather than ignored. Concretely, we stop treating a test result as a fixed ground truth and instead treat it as one sample from a per-chemical distribution, whose width the model learns and carries through to the downstream analysis.


Standard regression models, including the base Bayesian Factorization Machine of \citet{rendle2011bayesian}, assume a single global noise level shared by all observations. We consider this too restrictive for ecotoxicity, where some chemicals produce consistent results while others are intrinsically variable. Our contribution is to give each chemical its own noise term. Concretely, each observed value is modelled as the factorization machine mean plus Gaussian noise whose magnitude is specific to the chemical:
\begin{equation}
    y_i \mid \mathbf{x}_i, \theta, \alpha_{c_i} \sim \mathcal{N}\left( \hat{y}_i(\theta), \alpha_{c_i}^{-1} \right)
    \label{eq:likelihood}
\end{equation}
where $\alpha_c$ is the precision, the inverse of the noise variance, for chemical $c$. A high $\alpha_c$ describes a chemical whose tests are consistent, and a low $\alpha_c$ a chemical whose tests are highly variable. Because we keep all repeated measurements rather than averaging them, the model sees this spread directly and can estimate it.

The per-chemical noise terms are themselves tied together by a shared prior whose scale is learned from all chemicals at once. A chemical with many observations has its noise term set mainly by the scatter of its own measurements around the model mean. A chemical with few observations has its noise term pulled towards the level typical of the dataset as a whole, so that a handful of measurements does not produce an extreme or unstable estimate. This is the same partial-pooling idea applied to the noise, and it lets rarely tested chemicals borrow information from the rest.

The resulting per-chemical variance is the \emph{aleatoric} uncertainty, the irreducible measurement noise that more data cannot remove:
\begin{equation}
    \text{Var}_{\text{aleatoric}}(c^*) = \frac{1}{S} \sum_{s=1}^{S} \frac{1}{\alpha_{c^*}^{(s)}}
    \label{eq:aleatoric}
\end{equation}
It is a property of the chemical and does not vary across species or other features, a simplification we return to in the limitations. Together, the two sources combine into a total predictive variance for each prediction:
\begin{equation}
    \text{Var}_{\text{total}}(\mathbf{x}^*, c^*) = \text{Var}_{\text{epistemic}}(\mathbf{x}^*) + \text{Var}_{\text{aleatoric}}(c^*)
    \label{eq:total}
\end{equation}
For any species--chemical combination this total describes the range of LC50 values the model considers plausible, while keeping the two contributions separate: the epistemic part, which more testing could reduce, and the aleatoric part, which reflects the inherent noise of the test. Because the aleatoric part is a statement about the measurements themselves, it can be carried forward into hazard assessment, so that quantities derived from the predictions, such as Species Sensitivity Distributions and the hazardous concentrations taken from them, reflect both the model's uncertainty and the irreducible variability of the underlying tests. We show this in the Results.


\section{Results}

The results are presented in two parts. We first validate the raw outputs of the BFM, establishing that its point predictions are accurate and that its two uncertainty estimates are well-characterised. We then demonstrate how predictions and their associated uncertainty propagate to Species Sensitivity Distributions (SSDs) and the derivation of hazardous concentrations, one of the practical application forms of observed and predicted ecotoxicity data in environmental protection, assessment and management.

\subsection{Prediction accuracy}
The evaluation of the BFM predictive accuracy considers the comparison of observed and predicted LC50 values. Each predicted value is the model's best estimate of the LC50 for an untested species--chemical combination, a data-driven form of read-across whose use is explored in detail by \citet{PosthumaImprovingHazard}. The BFM model achieves an RMSE of 0.70 log mg/L in grouped 5-fold cross-validation ($k=32$, $n_{\text{iter}}=2000$, $n_{\text{burn}}=100$), where each held-out fold contains chemical--species--duration triplets never seen during training. This represents an improvement on \citet{PosthumaImprovingHazard}, who reported an RMSE of 0.85, the improvement likely attributable to the incorporation of chemical and species features in the model. Direct comparisons to other works are more difficult to make, for reasons described in  \citet{schur2025comparability}. \citet{Viljanen2023} tested a number of different models which yielded RMSE's in the 0.88-1 range, albeit on a different dataset. Achieving the best possible performance in this prediction setting is not the primary concern of this study, and therefore we do not attempt to further compare RMSE performance. 

The prediction accuracy of the BFM is illustrated in Figure~\ref{fig:pred_vs_observed}, which plots model predictions for species--chemical pairs at 48-hour exposure against the average of the actual observations for those pairs. The strong correlation gives a visual impression of how well the model recovers measured toxicity. Some scatter is expected from the observations themselves: each plotted x-value may rest on a single test or on many (Figure~\ref{fig:rsd-distr}), so the less-tested points carry more measurement noise. Consistent with this, \citet{PosthumaImprovingHazard} found that the predicted--observed correlation improves substantially when restricted to more-tested, robust datapoints.

\begin{figure}[htbp]
    \centering
    \includegraphics[width=0.6\linewidth]{images/predicted_vs_measured_48h.png}
    \caption{Each point on the plot is the mean measured log(mg/L) toxicity (average where multiple measurements available) vs.\ the mean predicted log(mg/L) toxicity for a unique chemical--species pair at 48 hours.}
    \label{fig:pred_vs_observed}
\end{figure}

A further inspection of the residuals in Figure~\ref{fig:residuals} allows for a deeper understanding of the biases of the model. Panel~(a) shows a scatterplot of the residuals against the predicted toxicity value. Aside from a few extremely outlying points, the shape of the plot is relatively uniform and suggests calibrated model behaviour. Panel~(b) shows the residuals scattered against the measured observations, showing a negative diagonal slope, indicating regression to the mean: very high measured values tend to be under-predicted while very low values tend to be over-predicted.
Having established that the model predicts accurately on average, the key question for risk assessment is whether we can trust individual predictions, particularly for chemicals with little testing history for the chemical--species pair itself.
\begin{figure}[htbp]
    \centering
    \includegraphics[width=1\linewidth]{images/bias_analysis_48h.png}
    \caption{Residual diagnostic plots showing (a) residuals vs predicted values, (b) residuals vs measured values revealing regression to the mean, and (c)  residual distribution histogram.}
    \label{fig:residuals}
\end{figure}


\subsection{Validating the two uncertainty components}

Each model prediction carries two uncertainty components. Epistemic uncertainty, derived from the variance across the model's posterior samples, is interpreted as the model confidence is in its mean prediction and it shrinks as more data constrains the fit. Aleatoric uncertainty is a per-chemical variance term the model learns directly, capturing the irreducible spread between repeat tests: the biological and experimental variability intrinsic to the assay, which represents the residual noise floor of the measurement rather than a gap in the model's features. Together these give a total predictive variance for each prediction; across all predictions, aleatoric variance usually contributes more than epistemic (${\sim}67\%$ of the total), 
indicating the model is relatively confident in its mean predictions but predicts spread in individual measurements, consistent with the high inter-test variability documented in Figure~\ref{fig:rsd-distr}. We validate these estimates in three steps: first that the aleatoric term recovers the spread seen in genuine repeat measurements, second that the two components respond to data availability as a reducible and an irreducible quantity should, and third that the combined predictive intervals are well-calibrated.

\subsubsection{Aleatoric uncertainty matches the spread of repeat measurements}

Repeated tests of the same species--chemical pair vary substantially (Figure~\ref{fig:rsd-distr}), so a single test result is best understood as one draw from a distribution rather than a fixed ground truth. The aleatoric term is the model's estimate of the width of that distribution, and we first check that it matches the spread of genuine repeat measurements.

Table~\ref{tab:replicate_stats} assesses whether the model's per-chemical aleatoric
uncertainty matches the empirical spread of genuine repeat measurements. For each
chemical--species--duration triplet measured at least five times, we compare the
empirical standard deviation of its replicates with the model's predicted aleatoric
standard deviation, summarising agreement by the ratio of the two. Across every
replicate bin the predicted and empirical standard deviations are of the same magnitude
(pooled predicted-to-empirical ratios range from $1.21$ at 5--9 replicates to $0.95$ at
100+), indicating that the learned aleatoric term recovers the scale of irreducible
measurement noise rather than absorbing systematic model error.

At low replicate counts the predicted spread is somewhat wider than the empirical one: at 5–9 replicates the bin-mean predicted SD exceeds the empirical SD by a pooled factor of about 1.2, with a median per-triplet ratio of 1.5. This gap closes as replicates accumulate — the median ratio falls monotonically to 0.90 by the 100+ bin and the pooled ratio to 0.95 — and it has a straightforward explanation that does not imply miscalibration. The comparison's denominator, the empirical SD, is itself computed from only a handful of replicates, and for small samples the sample standard deviation is both noisy and biased downward as an estimator of the true σ. The ratio is therefore inflated at low replicate counts even when the predicted SD is unbiased, and it settles toward one as the empirical estimate stabilises. The wide interquartile range in the sparse bins ([0.90, 2.69] at 5–9 replicates, tightening to [0.89, 1.03] at 100+) is a direct signature of this sampling noise in the denominator rather than of scatter in the model's estimate. A smaller, second contribution is the shared prior on the per-chemical precisions, which applies mild shrinkage to sparsely-measured chemicals toward the population's typical noise level.


For the data-rich chemicals that dominate regulatory SSD construction, the predicted and
empirical spreads coincide closely (Pooled Ratio $1.12$ at 50--99 replicates and $0.95$
at 100+), so the uncertainty estimates are most reliable in precisely the regime where
they are used. From a regulatory standpoint, the mild over-prediction that remains for
data-poor chemicals is conservative: it widens the confidence intervals on derived
hazard concentrations rather than producing false precision.

% TODO (LP Comment 343): Leo asks "Is this correctly worded? Is this what you mean here to be over-estimated?" 

% TODO (LP Comment 346): Leo says: "I would reserve the word 'conservative' for the
% regulatory-use context. What is the mathematical formulation here, solely driven by
% insights in the data and the predictions?" 
\begin{table}[ht]
    \centering
    \caption{Calibration of the model's aleatoric uncertainty against the empirical
    spread of genuine repeat measurements. For every unique chemical--species--duration
    triplet measured at least five times, the empirical standard deviation of the
    replicates is compared with the model's predicted aleatoric standard deviation;
    triplets are grouped into bins by replicate count. \textbf{N~groups}: number of
    triplets in the bin.  \textbf{Median Ratio} and \textbf{Ratio IQR}: median and
    interquartile range of the per-triplet predicted-to-empirical SD ratio. \textbf{Mean
    Pred~SD} and \textbf{Mean Emp~SD}: bin-mean predicted and empirical standard
    deviations (log mg/L). A ratio above one indicates the model predicts a wider spread
    than is observed empirically.}
    \label{tab:replicate_stats}
\begin{tabular}{lccccc}
    \toprule
    \textbf{Replicates} & \textbf{N groups} & \textbf{Median Ratio} & \textbf{Ratio IQR} & \textbf{Mean Pred SD} & \textbf{Mean Emp SD} \\
    \midrule
    5--9   & 1,803 & 1.474 & [0.90, 2.69] & 0.470 & 0.387 \\
    10--19 & 658   & 1.449 & [0.93, 2.47] & 0.467 & 0.381 \\
    20--49 & 227   & 1.315 & [0.93, 2.00] & 0.482 & 0.400 \\
    50--99 & 30    & 1.091 & [0.91, 1.51] & 0.520 & 0.463 \\
    100+   & 6     & 0.897 & [0.89, 1.03] & 0.533 & 0.558 \\
    \bottomrule
\end{tabular}
\end{table}

% TODO (LP Comment 385): Leo asks whether this figure (uncertainty vs observations)
% should be renumbered and positioned much earlier possibly before Table 2 or
% directly after the predictive accuracy subsection. 

\subsubsection{Epistemic and aleatoric uncertainty respond differently to data}

The two uncertainty components behave differently as data accrues, which is examined at the chemical level by plotting out-of-fold epistemic and aleatoric standard deviations against the total number of observations per chemical (Figure~\ref{fig:uncert_vs_num_obs}). Epistemic uncertainty showed a negative association with observation count (Spearman $\rho = -0.308$), consistent with the expectation that additional data constrains the chemical's latent factors and reduces parametric uncertainty: the reducible component shrinks where more evidence is available. This is the uncertainty that any statistical or model carries but that most methods leave unquantified; the posterior ensemble makes it explicit by measuring how much the predictions of different plausible model fits disagree.

The aleatoric component, by contrast, shows essentially no association with observation count (Spearman $\rho = 0.087$). This is what one expects of a genuinely irreducible noise term: the intrinsic biological and experimental variability of a chemical's toxicity tests should not systematically diminish simply because the chemical has been tested more often. The parameter $\alpha_c$ is estimated per chemical and captures the residual variance around the model's mean — measurement noise together with the inter‑species and inter‑duration spread that the rank‑32 factorisation does not resolve, and any remaining chemical‑specific misfit. Its flatness in observation count follows from the shared prior on the per‑chemical precisions. Because that prior's rate is learned from the whole population of chemicals, its scale tracks the typical per‑chemical residual variance in the data: a chemical with few observations has its precision pooled toward this population‑typical noise level, whereas one with many observations has its precision determined by the scatter of its own residuals. The pooling target is the same whether a chemical is sparsely or heavily tested, a chemical's intrinsic variability has no reason to depend on how often it happens to have been tested, so the estimate inherits no systematic trend with observation count. It is this pooling toward a data‑derived centre, rather than toward a fixed predetermined scale, that keeps the relationship flat. Epistemic and aleatoric uncertainty therefore behave exactly as a reducible and an irreducible quantity should: the first shrinks as evidence accumulates, the second stays broadly constant.

\begin{figure}[htbp]
    \centering
    \includegraphics[width=1\linewidth]{images/uncertainty_vs_observations_48h.png}
    \caption{Relationship between data availability and predicted uncertainty. Each point represents one chemical at 48\,h exposure duration. The x-axis shows the total number of observations (across all durations) available for that chemical in the training set. Chemicals with fewer than three total observations were excluded to avoid cold-start artifacts in the cross-validated uncertainty estimates (N\,=\,660 excluded; N\,=\,1,722 retained). (a) Mean epistemic (reducible) standard deviation (SD), computed as the average across out-of-fold posterior predictive variance for each chemical's 48\,h observations. (b) Mean aleatoric (non-reducible) SD, derived from the per-chemical precision parameter $\alpha_c$. Both uncertainty components are expressed in log mg/L units.}
    \label{fig:uncert_vs_num_obs}
\end{figure}

Figure~\ref{fig:variance_composition} summarises this contrast by decomposing the total predictive variance across observation-count bins. Epistemic variance declines steadily as observations accumulate, whereas aleatoric variance stays broadly stable; its apparent increase in the highest-count bin reflects the particular chemicals that bin contains rather than a systematic trend with data. Because epistemic shrinks while aleatoric does not, the share of total variance attributable to aleatoric noise grows with the amount of data: epistemic uncertainty dominates for sparsely tested chemicals, and aleatoric uncertainty dominates for well-tested ones.

% =========================================================================
% PLACEHOLDER: new variance-composition figure (the stacked-bar + share plot
% from discussion). Decide later whether to (a) include it here alongside
% Figure~\ref{fig:uncert_vs_num_obs}, or (b) replace Figure~\ref{fig:uncert_vs_num_obs}
% with it. If replacing, merge the two epistemic/aleatoric paragraphs above
% with the summary paragraph here. If dropping, delete the paragraph above too.
% =========================================================================
\begin{figure}[htbp]
    \centering
    \includegraphics[width=\linewidth]{images/variance_composition_vs_nobs.png}
    \caption{Composition of total predictive variance by data availability. Chemicals are binned by their number of observations. (a)~Mean predictive variance per bin, split into aleatoric (irreducible noise) and epistemic (model uncertainty) components.}
    \label{fig:variance_composition}
\end{figure}

The aleatoric parameter should not be interpreted as irreducible measurement noise in isolation: because $\alpha_c$ is shared across all species and durations for a given chemical, it absorbs any residual variance the model's rank-32 factorisation does not capture. The total predictive uncertainty for a given (chemical, species, duration) triplet is the sum of the epistemic and aleatoric components.

\subsubsection{Combined predictive intervals are well-calibrated}

To assess whether the total predictive uncertainty is well-calibrated, we construct posterior predictive intervals at a range of nominal coverage levels and measure how often the true held-out observation falls within them. For each out-of-fold observation and each posterior sample~$s$, we draw from $\mathcal{N}(\hat{y}^{(s)}, 1/\alpha_c^{(s)})$, combining the epistemic variation across posterior samples with the aleatoric noise learned for that chemical. Figure~\ref{fig:calibration} plots the resulting empirical coverage against the nominal level. The model is slightly conservative at low-to-mid levels (a 50\% credible interval captures approximately 57\% of observations) and close to nominal in the 90--95\% range most relevant for regulatory applications (87.7\% and 91.5\% empirical coverage, respectively). The mean absolute calibration error across all levels is 0.043. The model thus errs slightly toward overcoverage at mid-range levels and shows only mild undercoverage in the high-coverage range that matters most for regulatory use, with empirical coverage falling two to four percentage points below nominal at the 90--95\% levels.
%CHECK VALUES AFTER MODEL UPDATE OR AFTER FINAL RE-TRAINING

\begin{figure}[htbp]
    \centering
    \includegraphics[width=0.55\linewidth]{images/calibration_curve.png}
    \caption{Posterior predictive calibration evaluated on 70{,}670 out-of-fold observations. For each nominal credible interval level, empirical coverage reports the fraction of observations whose true value falls within the interval. The dashed line indicates perfect calibration.}
    \label{fig:calibration}
\end{figure}

To give a visual idea of what these uncertainty estimates look like in practice, Figure~\ref{fig:chosen_predictions} shows model predictions for five randomly selected chemicals alongside the species for which they were tested in the dataset, allowing side-by-side comparison of test results with predictions. The dark blue bars represent the epistemic model uncertainty around the mean, and the light blue bars extend to the total uncertainty, with the additional extent beyond the dark blue bars representing the per-chemical aleatoric spread.

% TODO (LP Comment 350): Leo notes that for some chemicals (e.g. Alachlor, a
% herbicide with specific MoA), observations may be dominated by algae tests that
% fit well, while predictions for other taxonomic groups may fail. He suggests
% showing chemicals grouped by MoA (narcotic, herbicide, insecticide) rather than
% random selection, and discussing "split-SSD" by taxonomic group.

% TODO (LP Comment 332): Leo suggests using the terms "data deserts" and "data
% mountains" to describe sparse vs dense regions of the heatmap, as an accessible
% metaphor for a lay audience.

\begin{figure}[htbp]
    \centering
    \includegraphics[width=0.85\linewidth]{images/example_predictions_48h.png}
    \caption{Predicted vs observed toxicity with decomposed uncertainty for five randomly selected chemicals at 48\,h exposure. For each chemical, the data on the tested species are sorted vertically by predicted toxicity. Blue dots show the model's mean prediction; dark blue bars indicate $\pm 1.96$ epistemic SD (posterior parametric uncertainty), and light blue bars extend to $\pm 1.96$ total SD (adding the per-chemical aleatoric component). Red crosses show the mean observed toxicity, with red bars indicating $\pm 1.96$ SD of replicate measurements where available.}
    \label{fig:chosen_predictions}
\end{figure}


\subsection{Propagating uncertainty into Species Sensitivity Distributions}

Species Sensitivity Distributions (SSDs) are one of the main ways ecotoxicity 
data are put to use in environmental protection, assessment, and management. An 
SSD describes how sensitivity to a chemical is distributed across species, and regulatory quantities are read off it as hazardous concentrations 
(HCx): the HC5 of a NOEC-based SSD underpins protective concentration standards, 
and the HC20 of an EC10-based SSD is used in life-cycle impact assessment. We use 
LC50 here as a methodological prototype, but the SSD machinery is the same for any 
endpoint. Traditionally, an SSD is obtained by fitting a lognormal curve to the 
few observed values available for a chemical. Building on the model-predicted SSD 
approach of \citet{PosthumaImprovingHazard}, which replaces that curve fit with 
an empirical CDF of model predictions covering every species in the dataset, we 
construct SSDs directly from the BFM's predictions. Crucially, each prediction now 
carries a characterised uncertainty, and the Bayesian structure of the BFM lets us 
propagate that uncertainty through the SSD and into the hazardous concentrations 
derived from it, which a point-prediction SSD cannot do. We present the results in 
two views: the uncertainty on each individual species prediction, and the 
uncertainty on the SSD curve as a whole.
% TODO (LP Comment 393): Leo asks for more explanation of the "ensemble SSD" concept
% upfront, defined clearly so that a layman would understand it intuitively. 

% TODO (LP Comment 394): Leo suggests reordering: show the per-point uncertainty
% view (current Figure 8 / \ref{fig:uncertainty_simple}) FIRST, to explain the raw
% BFM output per individual species prediction, and THEN introduce the ensemble SSD
% (current Figure 7 / \ref{fig:ssd_comparison}) as the aggregated end-user view.


The per-species uncertainty is shown directly by plotting individual species predictions with their decomposed uncertainty bars (Figure~\ref{fig:uncertainty_simple}). For each species, the dark blue bars show $\pm 1.96$ aleatoric standard deviations, the constant per-chemical noise term, and the light blue bars extend to $\pm 1.96$ total standard deviations, the additional width reflecting the species-specific epistemic component. This per-point view reveals where in the SSD the uncertainty concentrates and highlights the distinct behaviour of the two uncertainty components: the aleatoric term is constant across species within a chemical but is somewhat larger for the sparsely tested Chlorfenprop-methyl (SD~$\approx 0.64$) than for Atrazine (SD~$\approx 0.40$), reflecting modestly greater residual variability for the less-studied compound. The epistemic term varies across species and is also larger for Chlorfenprop-methyl, differing by roughly a third between the two chemicals. It is important to note that these per-point error bars describe uncertainty around individual predictions, not around the SSD curve as a whole; the ensemble approach below is needed for the latter.

\begin{figure}[htbp]
    \centering
    \includegraphics[width=0.495\linewidth]{images/atrazine_novel_uncertainty.png}
    \includegraphics[width=0.495\linewidth]{images/chlorfenprop_methyl_novel_uncertainty.png}
    \caption{Model-predicted SSDs for Atrazine (left) and Chlorfenprop-methyl (right) with per-species uncertainty bars. Dark blue: $\pm 1.96$ aleatoric SD (the constant per-chemical noise term); light blue: $\pm 1.96$ total SD (epistemic + aleatoric). The aleatoric component is constant within each chemical but somewhat larger for the sparsely tested Chlorfenprop-methyl.}
    \label{fig:uncertainty_simple}
\end{figure}

% TODO (LP Comment 395): Leo asks whether these SSD scatter dots could have shapes
% and colours for different taxonomic groups.

Figure~\ref{fig:ssd_comparison} shows the ensemble of predicted SSDs for two contrasting chemicals, each overlaid on the traditional lognormal SSD fitted to the observed species. For Atrazine, a data-rich compound with 54 tested species and a specific Mode of Action (photosynthesis inhibitor, herbicide), the posterior ensemble forms a narrow band. For Chlorfenprop-methyl (also a herbicide), which has observations for only 7 species, the ensemble fans out considerably, reflecting the greater epistemic uncertainty associated with a sparsely tested chemical. This contrast demonstrates the key practical value of the approach: the width of the ensemble directly communicates how much the SSD, and any regulatory quantity derived from it, should be trusted. For Chlorfenprop-methyl this width is driven by epistemic uncertainty, which is reducible: additional testing of this chemical would narrow the band. For a data-rich chemical the residual width instead reflects irreducible measurement noise that further testing would not remove.

% TODO (LP Comment 392): Leo suggests expanding this to a four-panel figure:
% (A) data-rich narcotic chemical, (B) data-poor narcotic, (C) data-rich herbicide,
% (D) data-poor herbicide. Dots coloured by taxonomic group (green = plants,
% blue = arthropods, red = vertebrates) to show that colours are mixed across the
% SSD for narcotics but grouped/clustered for specific-MoA chemicals.

For each of the $S = 1{,}900$ post-burn-in Gibbs samples retained during training, we evaluate the full factorisation machine equation at $\theta^{(s)}$ to obtain a deterministic toxicity prediction $\hat{y}_i^{(s)}$ for every species~$i$ for each chemical of interest. Each set of predictions is sorted to produce a complete SSD curve, yielding an ensemble of $S$ internally consistent SSDs. Because each curve is derived from a single coherent parameter draw, the correlated epistemic structure across species is preserved: the entire SSD shifts and rescales as the shared parameters vary, rather than individual points scattering independently. Credible intervals for any functional metric derived from an SSD, such as the HC5 or HC20 (Hazardous Concentration for 5\% or 20\% of the species, respectively), can then be obtained directly from the empirical distribution of that quantity across the posterior ensemble.

\begin{figure}[htbp]
     \centering
     \begin{subfigure}[b]{0.48\textwidth}
         \centering
         \includegraphics[width=\textwidth]{images/ssd_uncertainty_atrazine_48h.png}
         \caption{Atrazine (54 tested species)}
         \label{fig:ssd_atrazine}
     \end{subfigure}
     \hfill
     \begin{subfigure}[b]{0.48\textwidth}
         \centering
         \includegraphics[width=\textwidth]{images/ssd_uncertainty_chlorfenprop_methyl_48h.png}
         \caption{Chlorfenprop-methyl (7 tested species)}
         \label{fig:ssd_chlorfenprop}
     \end{subfigure}
     \caption{Ensemble SSDs for Atrazine (left) and Chlorfenprop-methyl (right) at 48\,h exposure, overlaid on the traditional lognormal SSD. The shaded band is the BFM 95\% credible region and the solid navy line its ensemble median, both constructed from 1{,}900 posterior samples; the red curve is the traditional lognormal SSD fitted to the observed species (black points). The ensemble is narrow for the data-rich Atrazine, where it closely tracks the traditional fit, and wide for the sparsely tested Chlorfenprop-methyl, where the traditional fit still falls within the credible band.}
     \label{fig:ssd_comparison}
\end{figure}

Overlaying the traditional lognormal SSD on the same axes shows where the two methods agree and how firmly that agreement is pinned down. For Atrazine the ensemble is narrow and the traditional fit sits almost on top of the ensemble median, the two hazard estimates agreeing closely (BFM HC20 $= -0.33$ against a traditional value of $-0.41$ log mg/L) and the traditional value falling well within the model's 95\% credible interval, $[-0.50, -0.17]$. For Chlorfenprop-methyl the ensemble instead spans more than a log unit: the traditional value ($-2.79$) still falls within the model's credible interval, $[-2.99, -1.55]$, but that interval is wide enough that any hazard concentration inferred from only seven tested species is left largely unconstrained. In both cases the traditional and model-predicted curves are read in the same way; what the ensemble adds is an explicit statement of how well the underlying data determine the SSD, and with it the hazardous concentration derived from it.


\subsection{Uncertainty-aware hazardous concentrations}

SSD models are used for two main purposes,  to derive Hazardous Concentrations that are used as protective standards (which is based on selecting a regulatory maximum-accepted Y-value, and deriving a concentration threshold as resulting X-value), and to derive the Potentially Affected Fraction of species (Y-value) given an ambient pollutant concentration (X-value). In this section, we illustrate how the ensemble SSD can be used for the former, to estimate an HCx value from the resulting SSD.

% TODO (LP Comment 409): Leo asks to expand the HC20 forest plot (Figure 10) with
% Y-axis information stating MoA and n_tests for each chemical (e.g.,
% "gamma-Cyhalothrin; Insecticide, n=xxx"). This would show (a) that data-rich
% chemicals have closer agreement between BFM and traditional, and (b) how MoA
% influences the pattern.

The ensemble of posterior SSDs enables direct derivation of hazardous concentrations with credible intervals. Figure~\ref{fig:hc20_30_chems} shows BFM-derived HC20 values with 95\% credible intervals alongside HC20 values obtained from traditional lognormal curve-fitting for a selection of chemicals. Two patterns are apparent: (i)~many of the traditionally derived HC20 values fall outside the BFM's 95\% credible interval, and (ii)~the BFM systematically predicts lower (that is, more protective) HC20 values.

This conservative bias is confirmed across the full set of 412 chemicals for which both approaches could be computed (Figure~\ref{fig:hc20correlation}; Pearson $r = 0.949$). The BFM-derived HC20 values fall consistently below the 1:1 line, particularly for chemicals with fewer observed species (the darker points in Figure~\ref{fig:hc20correlation}); the shift below the line grows as the number of tested species falls (Spearman $\rho = 0.29$ between species count and the signed HC20 difference). The relationship is a statistical tendency rather than a rule: individual chemicals fall on both sides of the 1:1 line, but the balance is strongly toward more protective BFM estimates. The likely explanation is that traditional SSDs are fitted to the small, often non-representative subset of species that have been tested in the laboratory, whereas the BFM predictions cover all 1{,}267 species in the dataset, drawing on cross-species and cross-chemical read-across. The broader species coverage may include taxa that are more sensitive but have simply never been tested against a given chemical. From a regulatory perspective, this conservatism means the BFM approach is less likely to underestimate hazard, a desirable property, since an SSD fitted to only a handful of species carries large sampling uncertainty and, by omitting untested and potentially more sensitive taxa, can overestimate the protective concentration.

Taken together, the SSD and HC results show that the BFM provides not only a hazard estimate for each chemical but also an explicit statement of how much evidence supports it, which is the information a risk assessor needs to decide how far to rely on a model-derived value.

\begin{figure}[htbp]
    \centering
    \includegraphics[width=0.8\linewidth]{images/hc20_forest_plot.png}
    \caption{HC20 values derived from BFM posterior ensemble (blue, with 95\% credible intervals) compared to HC20 values from traditional lognormal curve-fitting (red markers). Chemicals are ordered by BFM-derived HC20.}
    \label{fig:hc20_30_chems}
\end{figure}

\begin{figure}[htbp]
    \centering
    \includegraphics[width=0.65\linewidth]{images/hc20_correlation_trad_vs_mc.png}
    \caption{Correlation between traditional (lognormal-fitted) and BFM-derived HC20 values for 412 chemicals with at least 5 tested species (Pearson $r = 0.949$). Points are coloured by the number of observed species per chemical. The systematic shift below the 1:1 line indicates more protective BFM estimates, particularly for data-poor chemicals.}
    \label{fig:hc20correlation}
\end{figure}


\section{Conclusion}

This work builds on using model-predictions directly in hazard assessment of \citet{PosthumaImprovingHazard}, which replaces lognormal SSD curve-fitting with point predictions of LC50 for every species in the dataset. We extend that approach by quantifying the uncertainty attached to each prediction, making the predictions more trustworthy for practical use and giving a more faithful account of the variability in the testing data. The Hierarchical Bayesian Factorization Machine (BFM) provides, for every species--chemical pair, two distinct and interpretable uncertainty components: epistemic uncertainty, reflecting the model's lack of knowledge in under-explored regions of the chemical--species matrix, and aleatoric uncertainty, capturing the irreducible variability inherent to biological testing at the per-chemical level. This was achieved without loss of accuracy: the BFM reaches an RMSE of 0.70 log mg/L in grouped 5-fold cross-validation, improving on the 0.85 reported by \citet{PosthumaImprovingHazard}.

The two components behave as a trustworthy decomposition should. The hierarchical prior on the per-chemical precision parameter enables principled partial pooling: chemicals with abundant test data have their noise variance dominated by the empirical residuals, while data-poor chemicals are regularised toward the population mean. Epistemic uncertainty decreases with observation count, confirming that the model becomes more confident where more evidence exists, while aleatoric uncertainty shows no systematic trend with data, consistent with an irreducible noise term whose scale matches the spread of repeated laboratory tests. The combined intervals are well-calibrated at the coverage levels most relevant for regulatory use. Because the reducible component shrinks as data accrues while the irreducible component does not, the balance between the two depends on how well a chemical has been studied: predictions for sparsely tested chemicals are dominated by reducible model uncertainty, whereas predictions for well-studied chemicals are dominated by irreducible measurement noise.

We further show how this uncertainty propagates into the downstream use of ecotoxicity data for hazard assessment. Constructing SSDs directly from the ensemble of posterior samples yields SSDs with credible intervals that preserve the correlated epistemic structure across species, so the entire distribution shifts coherently with the shared parameters rather than individual points scattering independently. This enables uncertainty-aware derivation of Hazardous Concentrations such as the HC20, which correlate strongly with those from traditional lognormal curve-fitting while adding quantified credible intervals and covering all species in the dataset rather than only the tested subset. The width of these intervals communicates how much a derived hazard value should be trusted, and is widest for the data-poor chemicals where traditional small-sample SSDs are least reliable. Generating predictions with calibrated uncertainty for the 99.5\% of the chemical--species matrix that remains untested addresses a long-standing bottleneck in environmental risk assessment, equipping risk assessors with a distribution of plausible outcomes rather than a single point estimate, and supporting decisions that account for the evidence available for each chemical.

Two facts underlie this work: laboratory measurements are not perfect ground truth but noisy samples, and predictions from any statistical model carry uncertainty. Both can be modelled explicitly, and as shown here, doing so need not cost predictive accuracy. Neither the Bayesian reformulation that yields epistemic uncertainty nor the explicit per-chemical noise model that yields aleatoric uncertainty is specific to the factorization machine or to ecotoxicology. The same two steps can be applied to other models, and in other fields, wherever predictions are made from sparse, noisy, and unevenly sampled data and used to inform decisions, as recent uncertainty-aware work in adjacent areas has begun to do \citep{von2026uncertainty, iwasaki2026bayesian}. Treated this way, machine learning is not a black box: domain understanding and statistical modelling combine to show where a prediction can be trusted, where it cannot, and where the variability is an irreducible feature of the system that any assessment should take into account.

Finally, because the model reports where its predictions are least certain, it could guide the prioritisation of new tests towards the chemical--species combinations that would most reduce uncertainty. This is an uncertainty-guided strategy analogous to active learning, which could direct testing effort more efficiently and, over time, contribute to the goal of reducing animal testing.

\section{Limitations and future work}

Several limitations of this work should be noted, each indicating a direction for further development.

\begin{enumerate}
    \item Aleatoric uncertainty is modelled at the chemical level rather than the species or observation level. The precision parameter $\alpha_c$ is shared across all species and durations for a given chemical, so the model cannot represent the possibility that some species respond more variably to a given chemical than others. A more granular parameterisation could improve calibration but would increase the number of parameters and the data required to fit them. Resolving the aleatoric variance further, for example jointly across chemicals and species, is a natural next step.

    \item The model cannot extrapolate to entirely novel chemicals or species. The factorization machine relies on one-hot encoding of chemical and species identities, so predictions require both to have been observed at least once during training. Extending the model to untested chemicals would require richer molecular representations, such as molecular fingerprints or SMILES embeddings, together with phylogenetic features, and remains a major open challenge in the field.

    \item A mild over-prediction of aleatoric variance remains for data-poor chemicals. At low replicate counts the predicted standard deviation exceeds the empirical one by a pooled factor of roughly 1.2 (Table~\ref{tab:replicate_stats}), relaxing towards unity as replicate counts grow. As discussed in the results, this is largely explained by the downward bias of the empirical standard deviation when it is estimated from only a few replicates, rather than by miscalibration of the model. For regulatory use the residual effect is conservative, widening rather than narrowing the intervals on derived hazard concentrations.

    \item Only the LC50 mortality endpoint is considered. Incorporating additional endpoints, such as EC50 or NOEC, could improve predictive power by increasing the number of observations per chemical, but would require careful harmonisation of endpoint definitions.

    \item The model is predominantly data-driven, relying mainly on chemical and species identities together with a small set of features. This makes it sensitive to biases in the underlying dataset and limits its direct transfer to other settings. Applying the methodology to other public or proprietary datasets, and evaluating its applicability domain, would clarify how far the predictions and their uncertainties can be relied upon for new chemicals and species.
\end{enumerate}

% =====================================================================
% Supporting Information for
% "Uncertainty Quantification for Aquatic Ecotoxicity Prediction"
%
% Included as the final section of main.tex via % =====================================================================
% Supporting Information for
% "Uncertainty Quantification for Aquatic Ecotoxicity Prediction"
%
% Included as the final section of main.tex via % =====================================================================
% Supporting Information for
% "Uncertainty Quantification for Aquatic Ecotoxicity Prediction"
%
% Included as the final section of main.tex via \input{supporting_information}.
% Relies on the preamble of main.tex (amsmath, booktabs, etc. are already
% loaded there; amssymb is added there for \Real below). Uses the shared
% references.bib already cited by the rest of the paper.
% =====================================================================

% ---- convenience macros ----
\newcommand{\yhat}{\hat{y}}
\newcommand{\Norm}{\mathcal{N}}
\newcommand{\Gam}{\mathrm{Gamma}}
\newcommand{\Real}{\mathbb{R}}
\newcommand{\E}{\mathbb{E}}
\newcommand{\Var}{\mathrm{Var}}
\newcommand{\vv}{\mathbf{v}}
\newcommand{\xx}{\mathbf{x}}

\section{Supporting Information}
\label{sec:si}

This Supporting Information gives the full mathematical specification of the
Hierarchical Bayesian Factorization Machine (BFM) summarised in the main text.
It is written for a statistical or machine-learning audience and documents the
model exactly as implemented: the predictor, the complete hierarchical
generative model, the Gibbs sampler and every conditional posterior, the
empirical-Bayes initialisation, the hyperparameter settings, and the
posterior-predictive quantities used for calibration, Species Sensitivity
Distributions (SSDs), and hazardous concentrations. The predictive core is the
factorization machine (FM) of \citet{rendle2010fm} cast in the Bayesian form of
\citet{rendle2011bayesian}; Section~\ref{sec:diff} states precisely how our
model differs from theirs. Symbols follow the main text and are collected in
Table~\ref{tab:notation}.

% =====================================================================
\subsection{Notation}
\label{sec:notation}

\begin{table}[htbp]
\centering
\caption{Notation used throughout the Supporting Information.}
\label{tab:notation}
\small
\begin{tabular}{ll}
\toprule
\textbf{Symbol} & \textbf{Meaning} \\
\midrule
$n$ & number of observations (rows) \\
$p$ & number of features (design-matrix columns) \\
$k$ & number of latent factors per feature (factorization rank) \\
$C$ & number of chemicals (noise groups) \\
$\xx_i \in \Real^{p}$ & feature vector of observation $i$ \\
$x_{ij}$ & $j$-th entry of $\xx_i$ \\
$y_i$ & centred log-transformed LC50 of observation $i$ \\
$c_i \in \{1,\dots,C\}$ & chemical (group) index of observation $i$ \\
$w_0$ & global bias \\
$w_j$ & linear weight of feature $j$ \\
$v_{jf}$ & $f$-th latent factor of feature $j$; $\vv_j=(v_{j1},\dots,v_{jk})$ \\
$\alpha_c$ & noise precision of chemical $c$ (aleatoric term) \\
$\mu_w,\lambda_w$ & mean and precision of the shared prior on the weights $w_j$ \\
$\mu_v,\lambda_v$ & mean and precision of the shared prior on the factors $v_{jf}$ \\
$a_0,b_0$ & shape and rate of the shared prior on the $\alpha_c$ \\
$\mu_0,\gamma_0$ & mean and precision of the Gaussian mean-hyperpriors \\
$\alpha_\lambda,\beta_\lambda$ & shape and rate of the Gamma prior on $\lambda_w,\lambda_v$ \\
$c_0,d_0$ & shape and rate of the Gamma hyperprior on $b_0$ \\
$\theta$ & all model parameters $\{w_0,\{w_j\},\{v_{jf}\}\}$ \\
$S$ & number of retained (post-burn-in) posterior samples \\
$\yhat^{(s)}$ & prediction from posterior sample $s$ \\
\bottomrule
\end{tabular}
\end{table}

We write $\Norm(\mu,\sigma^2)$ for a Gaussian with mean $\mu$ and variance
$\sigma^2$, and $\Gam(a,b)$ for a Gamma distribution with shape $a$ and rate $b$
(so that its mean is $a/b$). All precisions are inverse variances.

% =====================================================================
\subsection{Data representation and feature encoding}
\label{sec:features}

Each observation is a single laboratory LC50 measurement for a
(chemical, species, duration) combination. The target $y_i$ is the
log-transformed concentration in mg/L, centred by subtracting the global mean
$\bar{y}$ of the training set; the mean is added back to recover log-mg/L units
when reporting predictions. Repeated measurements of the same triplet are kept
as separate rows rather than averaged, so that the per-chemical noise model
(Section~\ref{sec:model}) sees the within-triplet spread directly.

The feature vector $\xx_i$ concatenates five one-hot categorical blocks and two
standardised numerical features:
\begin{itemize}
  \item \textbf{chemical} identity (one-hot over $C$ chemicals),
  \item \textbf{species} identity (one-hot),
  \item \textbf{test duration} (one-hot over $\{24,48,72,96\}$\,h),
  \item \textbf{taxonomic family} and \textbf{taxonomic class} of the species (two one-hot blocks),
  \item \textbf{numerical}: log-transformed molecular weight and the RDKit
        octanol--water partition coefficient (\texttt{clogp}), median-imputed
        and standardised to zero mean and unit variance.
\end{itemize}
Unseen categories at prediction time map to an all-zero block (they contribute
nothing to the sums below), so the model can score any (chemical, species,
duration) triplet whose chemical and species were each observed at least once
during training, but not entirely novel chemicals or species. For the dataset
used here $n=70{,}670$, $C=3{,}295$ chemicals, $1{,}267$ species, and
$p=4{,}879$ (namely $3{,}295+1{,}267+4+283+28+2$).

% =====================================================================
\subsection{The factorization-machine predictor}
\label{sec:fm}

The predictive mean for observation $i$ is the second-order factorization
machine
\begin{equation}
  \yhat_i(\theta) \;=\; w_0 \;+\; \sum_{j=1}^{p} w_j\,x_{ij}
  \;+\; \sum_{j=1}^{p}\sum_{j'>j}^{p} \langle \vv_j,\vv_{j'}\rangle\,x_{ij}x_{ij'},
  \label{eq:fm}
\end{equation}
with $\langle \vv_j,\vv_{j'}\rangle=\sum_{f=1}^{k} v_{jf}v_{j'f}$. The pairwise
interactions are not stored explicitly; they are represented in the low-rank
form of Eq.~\eqref{eq:fm}, which lets the model infer the interaction between a
chemical and a species that were never tested together. Using the standard FM
identity, the interaction sum is evaluated in $O(pk)$ time (in practice $O(k
N_{\mathrm{nz}})$ for a sparse $\xx_i$ with $N_{\mathrm{nz}}$ non-zeros) as
\begin{equation}
  \sum_{j<j'} \langle \vv_j,\vv_{j'}\rangle\,x_{ij}x_{ij'}
  \;=\; \tfrac{1}{2}\sum_{f=1}^{k}\!\left[
      \Big(\sum_{j=1}^{p} v_{jf}x_{ij}\Big)^{\!2}
      - \sum_{j=1}^{p} v_{jf}^{2}x_{ij}^{2}
  \right].
  \label{eq:fm-identity}
\end{equation}
The diagonal subtraction in Eq.~\eqref{eq:fm-identity} removes each feature's
self-interaction. A direct consequence, used repeatedly by the sampler, is that
$\yhat_i$ is an \emph{affine} function of each individual parameter when the
others are held fixed (Section~\ref{sec:sampler}).

% =====================================================================
\subsection{Full hierarchical generative model}
\label{sec:model}

The model is a Bayesian FM in which the single global noise precision of the
standard BFM is replaced by a per-chemical precision $\alpha_c$ under a shared
Gamma prior. The complete generative model is:

\paragraph{Likelihood.} Each observation is its FM mean plus Gaussian noise
whose precision is that of the observation's chemical,
\begin{equation}
  y_i \;\mid\; \xx_i,\theta,\alpha_{c_i}
  \;\sim\; \Norm\!\big(\yhat_i(\theta),\ \alpha_{c_i}^{-1}\big),
  \qquad i=1,\dots,n.
  \label{eq:like}
\end{equation}

\paragraph{Priors on the model parameters.} The global bias has a fixed Gaussian
prior; the weights and factors are drawn from shared populations whose mean and
precision are themselves learned:
\begin{align}
  w_0     &\sim \Norm(\mu_0,\ \gamma_0^{-1}), \\
  w_j     &\sim \Norm(\mu_w,\ \lambda_w^{-1}), & j&=1,\dots,p, \label{eq:prior-w}\\
  v_{jf}  &\sim \Norm(\mu_v,\ \lambda_v^{-1}), & j&=1,\dots,p,\ f=1,\dots,k. \label{eq:prior-v}
\end{align}

\paragraph{Hyperpriors on the populations.} There is one population for all
linear weights and one for all latent factors (as in the base BFM of
\citet{rendle2011bayesian}); each has a Gaussian hyperprior on its mean and a
Gamma hyperprior on its precision:
\begin{align}
  \mu_w &\sim \Norm(\mu_0,\ \gamma_0^{-1}), & \lambda_w &\sim \Gam(\alpha_\lambda,\ \beta_\lambda), \\
  \mu_v &\sim \Norm(\mu_0,\ \gamma_0^{-1}), & \lambda_v &\sim \Gam(\alpha_\lambda,\ \beta_\lambda).
\end{align}

\paragraph{Hierarchical noise model.} This is the methodological contribution.
Each chemical has its own noise precision, drawn from a Gamma prior whose rate is
shared across all chemicals and is itself given a Gamma hyperprior:
\begin{align}
  \alpha_c &\sim \Gam(a_0,\ b_0), & c&=1,\dots,C, \label{eq:prior-alpha}\\
  b_0      &\sim \Gam(c_0,\ d_0). \label{eq:prior-b0}
\end{align}
The shape $a_0$ is held fixed (set by empirical Bayes, Section~\ref{sec:eb});
the rate $b_0$ is resampled every sweep. Because $b_0$ is common to all
chemicals, it ties the per-chemical precisions together: a chemical with many
observations has its $\alpha_c$ set mainly by the scatter of its own residuals,
whereas a sparsely measured chemical is pooled towards the population-typical
noise level implied by $\Gam(a_0,b_0)$. This partial pooling is what keeps the
aleatoric estimate stable for data-poor chemicals.

The joint density factorises as
\begin{equation}
\begin{aligned}
  p\big(\theta,\{\alpha_c\},b_0,\mu_w,\lambda_w,\mu_v,\lambda_v \mid \mathbf{y}\big)
  \;\propto\;
  & \Big[\textstyle\prod_{i=1}^{n} \Norm\!\big(y_i \mid \yhat_i(\theta),\alpha_{c_i}^{-1}\big)\Big]
    \Norm(w_0\mid\mu_0,\gamma_0^{-1}) \\
  &\times \Big[\textstyle\prod_{j=1}^{p}\Norm(w_j\mid\mu_w,\lambda_w^{-1})\Big]
          \Big[\textstyle\prod_{j,f}\Norm(v_{jf}\mid\mu_v,\lambda_v^{-1})\Big] \\
  &\times \Norm(\mu_w\mid\mu_0,\gamma_0^{-1})\,\Gam(\lambda_w\mid\alpha_\lambda,\beta_\lambda)\,
          \Norm(\mu_v\mid\mu_0,\gamma_0^{-1})\,\Gam(\lambda_v\mid\alpha_\lambda,\beta_\lambda) \\
  &\times \Big[\textstyle\prod_{c=1}^{C}\Gam(\alpha_c\mid a_0,b_0)\Big]\,\Gam(b_0\mid c_0,d_0).
\end{aligned}
\label{eq:joint}
\end{equation}
This is the model drawn in plate notation in the main text: the fixed constants
$(\mu_0,\gamma_0,\alpha_\lambda,\beta_\lambda,a_0,c_0,d_0)$ set the hyperpriors;
$w_j$ and $v_{jf}$ share $(\mu_w,\lambda_w)$ and $(\mu_v,\lambda_v)$
respectively; the $\alpha_c$ share the rate $b_0$; and each $y_i$ is drawn from a
Gaussian centred at $\yhat_i$ with precision $\alpha_{c_i}$.

% =====================================================================
\subsection{Relationship to the standard Bayesian FM}
\label{sec:diff}

Our model departs from the BFM of \citet{rendle2011bayesian} in two respects.
\begin{enumerate}
  \item \textbf{Per-chemical noise (the substantive change).} The standard BFM
  assumes a single global precision $\alpha$ shared by all observations, with a
  Gamma prior and a Gamma-conjugate posterior. We replace it by $C$
  chemical-specific precisions $\alpha_c$ (Eq.~\ref{eq:prior-alpha}) under a
  shared, learned rate $b_0$ (Eq.~\ref{eq:prior-b0}). Equivalently, the global
  $\alpha$ in every conditional below becomes the observation-specific
  $\alpha_{c_i}$, and one extra conjugate draw for $b_0$ is added per sweep. This
  turns the noise term into a per-chemical, partially pooled quantity, which is
  the aleatoric component carried into the downstream analysis.

  \item \textbf{Independent Gaussian priors on the population means (a minor
  change).} The standard BFM couples each population mean to its precision
  through a Normal--Gamma prior, $\mu \mid \lambda \sim
  \Norm(\mu_0,(\gamma_0\lambda)^{-1})$, so the posterior precision of $\mu$
  scales with $\lambda$. We instead place an independent
  $\Norm(\mu_0,\gamma_0^{-1})$ prior on each mean (Eqs.~\ref{eq:prior-w}--\ref{eq:prior-v}),
  decoupling $\mu$ from $\lambda$; the mean does not then enter the update for
  $\lambda$, and the $\mu$-posterior precision is $\gamma_0+p\lambda$ rather than
  $(\gamma_0+p)\lambda$ (Section~\ref{sec:sampler}). With $p$ in the thousands
  this has negligible effect on inference, as \citet{rendle2011bayesian} also
  note for the hyperprior constants, but we state it for exactness.
\end{enumerate}
As in the reference, we use single-parameter (unblocked) Gibbs sampling, which
avoids the matrix inversions of blocked schemes and keeps each sweep linear in
$k$.

% =====================================================================
\subsection{Gibbs sampler and conditional posteriors}
\label{sec:sampler}

Posterior inference is by Gibbs sampling: each quantity is drawn in turn from its
full conditional given the current values of the others. All conditionals are
available in closed form because (i) the FM mean is affine in each single
parameter, giving Gaussian conditionals with Gaussian priors, and (ii) the
Gamma priors on the precisions are conjugate to the Gaussian likelihood.

Throughout we cache the running residual
\begin{equation}
  e_i \;=\; y_i - \yhat_i(\theta),
  \label{eq:resid}
\end{equation}
and the per-factor cache $q_{if}=\sum_{j=1}^{p} v_{jf}x_{ij}$, both updated by a
rank-one correction after each parameter is resampled. This is what makes a full
sweep $O(k N_{\mathrm{nz}})$, where $N_{\mathrm{nz}}$ is the number of non-zeros
in the design matrix. For a parameter with current value $\theta_\ell$ and
``slope'' $h_{i\ell}=\partial \yhat_i/\partial\theta_\ell$ (constant in
$\theta_\ell$ by multilinearity), the partial residual that adds back its current
contribution is $r_{i}=e_i+h_{i\ell}\theta_\ell$.

\subsubsection{Global bias $w_0$}
Here $h_{i0}=1$ for all $i$, so with prior $\Norm(\mu_0,\gamma_0^{-1})$,
\begin{equation}
  \sigma_{w_0}^2=\Big(\gamma_0+\sum_{i=1}^{n}\alpha_{c_i}\Big)^{-1},
  \qquad
  \mu_{w_0}=\sigma_{w_0}^2\Big(\gamma_0\mu_0+\sum_{i=1}^{n}\alpha_{c_i}\,(e_i+w_0)\Big),
  \qquad
  w_0 \sim \Norm(\mu_{w_0},\sigma_{w_0}^2).
\end{equation}
Then set $e_i \mathrel{+}= (w_0^{\mathrm{old}}-w_0^{\mathrm{new}})$ for all $i$.

\subsubsection{Weight population $(\mu_w,\lambda_w)$}
Treating the $p$ weights as i.i.d.\ $\Norm(\mu_w,\lambda_w^{-1})$ draws,
\begin{align}
  \mu_w &\sim \Norm\!\big(\mu_{\mu_w},\,\sigma_{\mu_w}^2\big),
  & \sigma_{\mu_w}^2 &=\big(\gamma_0+p\,\lambda_w\big)^{-1},
  & \mu_{\mu_w} &=\sigma_{\mu_w}^2\big(\gamma_0\mu_0+\lambda_w\textstyle\sum_{j=1}^{p}w_j\big), \\
  \lambda_w &\sim \Gam\!\Big(\alpha_\lambda+\tfrac{p}{2},\ \
     \beta_\lambda+\tfrac{1}{2}\textstyle\sum_{j=1}^{p}(w_j-\mu_w)^2\Big).
\end{align}

\subsubsection{Linear weights $w_j$}
For feature $j$, $h_{ij}=x_{ij}$, and only observations with $x_{ij}\neq 0$
contribute. With prior $\Norm(\mu_w,\lambda_w^{-1})$ and partial residual
$r_i=e_i+x_{ij}w_j$,
\begin{equation}
  \sigma_{w_j}^2=\Big(\lambda_w+\sum_{i:\,x_{ij}\neq0}\!\alpha_{c_i}\,x_{ij}^2\Big)^{-1},
  \quad
  \mu_{w_j}=\sigma_{w_j}^2\Big(\lambda_w\mu_w+\sum_{i:\,x_{ij}\neq0}\!\alpha_{c_i}\,x_{ij}\,r_i\Big),
  \quad
  w_j\sim\Norm(\mu_{w_j},\sigma_{w_j}^2).
\end{equation}
Then update $e_i \mathrel{+}= x_{ij}\,(w_j^{\mathrm{old}}-w_j^{\mathrm{new}})$.

\subsubsection{Factor population $(\mu_v,\lambda_v)$}
Treating all $pk$ factor entries as i.i.d.\ $\Norm(\mu_v,\lambda_v^{-1})$,
\begin{align}
  \mu_v &\sim \Norm\!\big(\mu_{\mu_v},\,\sigma_{\mu_v}^2\big),
  & \sigma_{\mu_v}^2 &=\big(\gamma_0+p\,k\,\lambda_v\big)^{-1},
  & \mu_{\mu_v} &=\sigma_{\mu_v}^2\big(\gamma_0\mu_0+\lambda_v\textstyle\sum_{j,f}v_{jf}\big), \\
  \lambda_v &\sim \Gam\!\Big(\alpha_\lambda+\tfrac{p k}{2},\ \
     \beta_\lambda+\tfrac{1}{2}\textstyle\sum_{j,f}(v_{jf}-\mu_v)^2\Big).
\end{align}

\subsubsection{Latent factors $v_{jf}$}
By the multilinearity of Eq.~\eqref{eq:fm-identity}, $\yhat_i$ is affine in
$v_{jf}$ with slope
\begin{equation}
  h_{ijf}=x_{ij}\big(q_{if}-x_{ij}v_{jf}\big),
  \qquad q_{if}=\sum_{j'=1}^{p} v_{j'f}x_{ij'},
\end{equation}
which does not depend on $v_{jf}$ (the $x_{ij}v_{jf}$ subtraction removes it).
With prior $\Norm(\mu_v,\lambda_v^{-1})$ and partial residual
$r_i=e_i+h_{ijf}v_{jf}$,
\begin{equation}
  \sigma_{v_{jf}}^2=\Big(\lambda_v+\sum_{i:\,x_{ij}\neq0}\!\alpha_{c_i}\,h_{ijf}^2\Big)^{-1},
  \quad
  \mu_{v_{jf}}=\sigma_{v_{jf}}^2\Big(\lambda_v\mu_v+\sum_{i:\,x_{ij}\neq0}\!\alpha_{c_i}\,h_{ijf}\,r_i\Big),
  \quad
  v_{jf}\sim\Norm(\mu_{v_{jf}},\sigma_{v_{jf}}^2).
\end{equation}
Then update both caches: $e_i \mathrel{+}= h_{ijf}\,(v_{jf}^{\mathrm{old}}-v_{jf}^{\mathrm{new}})$
and $q_{if}\mathrel{+}= x_{ij}\,(v_{jf}^{\mathrm{new}}-v_{jf}^{\mathrm{old}})$.

\subsubsection{Per-chemical noise precisions $\alpha_c$}
Let $N_c=|\{i:c_i=c\}|$ and $\mathrm{SSE}_c=\sum_{i:\,c_i=c} e_i^2$ be the
observation count and residual sum of squares of chemical $c$. Conjugacy of the
Gamma prior (Eq.~\ref{eq:prior-alpha}) to the Gaussian likelihood
(Eq.~\ref{eq:like}) gives
\begin{equation}
  \alpha_c \;\sim\; \Gam\!\Big(a_0+\tfrac{N_c}{2},\ \ b_0+\tfrac{1}{2}\,\mathrm{SSE}_c\Big),
  \qquad c=1,\dots,C.
  \label{eq:post-alpha}
\end{equation}
All $C$ precisions are drawn in one vectorised step using per-group sums of
$e_i^2$.

\subsubsection{Shared noise rate $b_0$}
The rate $b_0$ is conjugate to the $C$ Gamma-distributed precisions with common
shape $a_0$: combining $\prod_c \Gam(\alpha_c\mid a_0,b_0)\propto
b_0^{C a_0}e^{-b_0\sum_c\alpha_c}$ with the prior $\Gam(c_0,d_0)$ yields
\begin{equation}
  b_0 \;\sim\; \Gam\!\Big(c_0+C\,a_0,\ \ d_0+\textstyle\sum_{c=1}^{C}\alpha_c\Big).
  \label{eq:post-b0}
\end{equation}

\subsubsection{One sweep and sample collection}
A single Gibbs sweep applies, in order: $w_0$; $(\mu_w,\lambda_w)$ then every
$w_j$; $(\mu_v,\lambda_v)$ then every $v_{jf}$; the $\alpha_c$
(Eq.~\ref{eq:post-alpha}); and $b_0$ (Eq.~\ref{eq:post-b0}). The chain is run
for $n_{\mathrm{iter}}$ sweeps from an arbitrary initialisation
($w_0=0$, $w=\mathbf{0}$, $v_{jf}\sim\Norm(0,0.1^2)$, $\alpha_c=1$); the first
$n_{\mathrm{burn}}$ sweeps are discarded as burn-in and the remaining
$S=n_{\mathrm{iter}}-n_{\mathrm{burn}}$ sweeps are retained. Each retained sample
$s$ stores $\{w_0^{(s)},w^{(s)},v^{(s)},\{\alpha_c^{(s)}\},b_0^{(s)}\}$ and
corresponds to one complete, internally consistent parameterisation of the
model.

% =====================================================================
\subsection{Empirical-Bayes initialisation of the prior shape $a_0$}
\label{sec:eb}

The shared shape $a_0$ controls how strongly sparsely tested chemicals are
pooled towards the population noise level. Rather than fix it arbitrarily, we set
it by a method-of-moments empirical-Bayes pre-pass so that the prior matches the
dispersion of \emph{well-estimated} precisions:
\begin{enumerate}
  \item Run a short Gibbs pre-pass on the full training set with a non-informative,
  fixed noise prior ($a_0=1$, $b_0=1$, $b_0$ not learned) and collect the
  posterior draws $\{\alpha_c^{(s)}\}$.
  \item Restrict to data-rich chemicals ($N_c\ge N_{\min}$, with $N_{\min}=50$),
  whose posterior precisions are essentially data-driven, and form their
  posterior means $\hat{\alpha}_c=\tfrac{1}{S}\sum_s\alpha_c^{(s)}$.
  \item Fit a Gamma to $\{\hat{\alpha}_c\}$ by matching mean $m$ and variance
  $\nu$, and take the shape
  \begin{equation}
    a_0 \;=\; m^2/\nu.
  \end{equation}
  The corresponding rate is discarded, because the main run learns $b_0$ from its
  hyperprior (Eq.~\ref{eq:post-b0}).
\end{enumerate}
For the run reported in the paper this gives $a_0\approx1.75$; the main-run
posterior for the shared rate concentrates near $b_0\approx0.08$, so
$\Gam(a_0,b_0)$ is a weakly informative population prior on the per-chemical
precisions. The same EB $a_0$ is used for the cross-validation runs and for the
full-data model used to generate SSDs.

% =====================================================================
\subsection{Hyperparameter settings}
\label{sec:hyper}

Table~\ref{tab:hyper} lists the fixed constants and run settings. The mean- and
weight/factor-precision hyperpriors are the same trivial constants used by
\citet{rendle2011bayesian}; with $p$ in the thousands their influence on the
posterior is negligible. The noise-model hyperprior $\Gam(c_0,d_0)$ on $b_0$ and
the EB shape $a_0$ are specific to this work.

\begin{table}[htbp]
\centering
\caption{Hyperparameter and run settings.}
\label{tab:hyper}
\small
\begin{tabular}{lll}
\toprule
\textbf{Quantity} & \textbf{Symbol} & \textbf{Value} \\
\midrule
Factorization rank            & $k$                       & $32$ \\
Mean-hyperprior mean          & $\mu_0$                   & $0$ \\
Mean-hyperprior precision     & $\gamma_0$                & $1$ \\
Weight/factor precision prior & $\alpha_\lambda,\beta_\lambda$ & $1,\ 1$ \\
Noise-rate hyperprior         & $c_0,\ d_0$               & $1,\ 1$ \\
Noise-prior shape (empirical Bayes) & $a_0$               & $\approx 1.75$ (fixed after EB) \\
Noise-prior rate              & $b_0$                     & initialised at $1$, learned \\
Main run                      & $n_{\mathrm{iter}},\ n_{\mathrm{burn}}$ & $2000,\ 100$ \\
Retained samples              & $S$                       & $1900$ \\
EB pre-pass                   & $n_{\mathrm{iter}},\ n_{\mathrm{burn}}$ & $100,\ 50$ \\
EB data-rich threshold        & $N_{\min}$                & $50$ \\
Cross-validation              & folds                     & $5$ (grouped by triplet) \\
\bottomrule
\end{tabular}
\end{table}

% =====================================================================
\subsection{Posterior-predictive quantities}
\label{sec:predictive}

For a query $\xx^\ast$ belonging to chemical $c^\ast$, each retained sample gives
a deterministic FM prediction $\yhat^{(s)}(\xx^\ast)$ via Eq.~\eqref{eq:fm}. The
point prediction is the sample mean
\begin{equation}
  \yhat^\ast=\frac{1}{S}\sum_{s=1}^{S}\yhat^{(s)}(\xx^\ast).
\end{equation}
The predictive variance decomposes by the law of total variance,
$\Var(y^\ast\mid\mathcal{D})=\Var_\theta\!\big(\E[y^\ast\mid\theta]\big)
+\E_\theta\!\big[\Var(y^\ast\mid\theta)\big]$, whose two Monte-Carlo estimates are
the epistemic and aleatoric components:
\begin{align}
  \Var_{\mathrm{epi}}(\xx^\ast) &= \frac{1}{S}\sum_{s=1}^{S}\big(\yhat^{(s)}(\xx^\ast)-\yhat^\ast\big)^2,
  \label{eq:epi}\\
  \Var_{\mathrm{ale}}(c^\ast) &= \frac{1}{S}\sum_{s=1}^{S}\frac{1}{\alpha_{c^\ast}^{(s)}},
  \label{eq:ale}\\
  \Var_{\mathrm{tot}}(\xx^\ast,c^\ast) &= \Var_{\mathrm{epi}}(\xx^\ast)+\Var_{\mathrm{ale}}(c^\ast).
  \label{eq:tot}
\end{align}
The epistemic term is the spread of the mean across plausible fits and shrinks as
data constrain the parameters; the aleatoric term is the posterior-mean noise
variance of the chemical and does not vanish with more data. The aleatoric term
depends only on the chemical, so it is constant across species and durations for
a given chemical.

\paragraph{Calibration.} To assess whether $\Var_{\mathrm{tot}}$ is well
calibrated we form posterior-predictive draws that combine both components. For
each held-out observation and each sample $s$ we draw
\begin{equation}
  \tilde{y}^{(s)} \sim \Norm\!\big(\yhat^{(s)}(\xx^\ast),\ 1/\alpha_{c^\ast}^{(s)}\big),
\end{equation}
giving an $S$-point predictive sample. For a nominal central level $\rho$ we take
the empirical $[(1-\rho)/2,\ (1+\rho)/2]$ quantiles of these draws as the
credible interval, and report empirical coverage as the fraction of held-out
observations whose true value lies inside. Coverage is evaluated on the
out-of-fold observations (Section~\ref{sec:cv}), so predictions are scored by
models that did not train on them.

% =====================================================================
\subsection{SSD ensembles and hazardous concentrations}
\label{sec:ssd}

The Bayesian structure lets uncertainty propagate coherently into Species
Sensitivity Distributions. For a target chemical $c^\ast$ at a fixed duration,
and each retained sample $s$, we evaluate Eq.~\eqref{eq:fm} at $\theta^{(s)}$ for
all species $i=1,\dots,M$ and sort the resulting predictions ascending to obtain
one SSD curve $\big(\yhat^{(s)}_{(1)}\le\cdots\le\yhat^{(s)}_{(M)}\big)$ plotted
against the plotting positions $m/(M+1)$. Because every curve comes from a single
coherent parameter draw, the correlated epistemic structure across species is
preserved: the whole curve shifts and rescales as the shared parameters vary,
rather than species scattering independently. The ensemble of $S$ curves is
summarised pointwise by its median and central $95\%$ band.

A hazardous concentration $\mathrm{HC}x$ (the concentration protective of
$100-x\%$ of species) is read off each sampled curve at rank
$\lfloor x\,M/100\rfloor$, giving a posterior sample of $\mathrm{HC}x$ values;
its median and $2.5/97.5$ percentiles are the reported estimate and $95\%$
credible interval. Optionally, aleatoric noise can be injected per species by
adding a draw $\Norm\!\big(0,1/\alpha_{c^\ast}^{(s)}\big)$ to each prediction
before sorting; the SSD figures and HC intervals in the main text use the
epistemic ensemble (no added aleatoric noise) unless stated otherwise.

% =====================================================================
\subsection{Cross-validation and out-of-fold uncertainty}
\label{sec:cv}

Predictive accuracy and calibration are evaluated with grouped $5$-fold
cross-validation. Folds are formed by \emph{grouped} splitting on the
(chemical, species, duration) triplet identifier, so that all replicate
measurements of a held-out triplet are removed from training together and no
replicate of a scored triplet is ever seen during training. Each fold refits the
full model (including the hierarchical noise model, with the EB shape $a_0$ fixed
and $b_0$ learned) on its training partition and predicts the held-out partition;
concatenating the held-out predictions over folds yields the out-of-fold arrays
used for the RMSE, the uncertainty-versus-data analyses
(Eqs.~\ref{eq:epi}--\ref{eq:ale}), and the calibration curve. The full-data model
used to generate the SSDs and the all-chemical HC comparison is fit once on all
$n$ observations with the same settings.

Two practical points. First, the epistemic estimate of Eq.~\eqref{eq:epi} is only
meaningful for chemicals that appear in the training partition of a fold;
chemicals with fewer total observations than the number of folds can have folds
in which none of their rows are present, leaving their chemical-specific
parameters at the prior and deflating their apparent epistemic variance. Such
cold-start chemicals are excluded from the per-chemical
uncertainty-versus-data analyses in the main text. Second, the aleatoric
out-of-fold value for an observation is the sample mean of
$1/\alpha_{c}^{(s)}$ over that fold's retained samples, so it reflects the noise
precision learned from the fold's training data only.

% =====================================================================
\subsection{Computational complexity}
\label{sec:complexity}

Every conditional in Section~\ref{sec:sampler} is a single Gaussian or Gamma
draw, and the residual and factor caches make each parameter update a rank-one
correction touching only the non-zeros of the relevant feature column. One full
sweep therefore costs $O(k\,N_{\mathrm{nz}})$, linear in the factorization rank
and in the number of non-zero design-matrix entries, with an additional
$O(C)$ for the noise block (Eqs.~\ref{eq:post-alpha}--\ref{eq:post-b0}). This
matches the complexity of the unblocked BFM sampler of
\citet{rendle2011bayesian}; the per-chemical noise model adds only the vectorised
$\alpha_c$ draw and one scalar $b_0$ draw per sweep, so the hierarchical
extension does not change the asymptotic cost.

.
% Relies on the preamble of main.tex (amsmath, booktabs, etc. are already
% loaded there; amssymb is added there for \Real below). Uses the shared
% references.bib already cited by the rest of the paper.
% =====================================================================

% ---- convenience macros ----
\newcommand{\yhat}{\hat{y}}
\newcommand{\Norm}{\mathcal{N}}
\newcommand{\Gam}{\mathrm{Gamma}}
\newcommand{\Real}{\mathbb{R}}
\newcommand{\E}{\mathbb{E}}
\newcommand{\Var}{\mathrm{Var}}
\newcommand{\vv}{\mathbf{v}}
\newcommand{\xx}{\mathbf{x}}

\section{Supporting Information}
\label{sec:si}

This Supporting Information gives the full mathematical specification of the
Hierarchical Bayesian Factorization Machine (BFM) summarised in the main text.
It is written for a statistical or machine-learning audience and documents the
model exactly as implemented: the predictor, the complete hierarchical
generative model, the Gibbs sampler and every conditional posterior, the
empirical-Bayes initialisation, the hyperparameter settings, and the
posterior-predictive quantities used for calibration, Species Sensitivity
Distributions (SSDs), and hazardous concentrations. The predictive core is the
factorization machine (FM) of \citet{rendle2010fm} cast in the Bayesian form of
\citet{rendle2011bayesian}; Section~\ref{sec:diff} states precisely how our
model differs from theirs. Symbols follow the main text and are collected in
Table~\ref{tab:notation}.

% =====================================================================
\subsection{Notation}
\label{sec:notation}

\begin{table}[htbp]
\centering
\caption{Notation used throughout the Supporting Information.}
\label{tab:notation}
\small
\begin{tabular}{ll}
\toprule
\textbf{Symbol} & \textbf{Meaning} \\
\midrule
$n$ & number of observations (rows) \\
$p$ & number of features (design-matrix columns) \\
$k$ & number of latent factors per feature (factorization rank) \\
$C$ & number of chemicals (noise groups) \\
$\xx_i \in \Real^{p}$ & feature vector of observation $i$ \\
$x_{ij}$ & $j$-th entry of $\xx_i$ \\
$y_i$ & centred log-transformed LC50 of observation $i$ \\
$c_i \in \{1,\dots,C\}$ & chemical (group) index of observation $i$ \\
$w_0$ & global bias \\
$w_j$ & linear weight of feature $j$ \\
$v_{jf}$ & $f$-th latent factor of feature $j$; $\vv_j=(v_{j1},\dots,v_{jk})$ \\
$\alpha_c$ & noise precision of chemical $c$ (aleatoric term) \\
$\mu_w,\lambda_w$ & mean and precision of the shared prior on the weights $w_j$ \\
$\mu_v,\lambda_v$ & mean and precision of the shared prior on the factors $v_{jf}$ \\
$a_0,b_0$ & shape and rate of the shared prior on the $\alpha_c$ \\
$\mu_0,\gamma_0$ & mean and precision of the Gaussian mean-hyperpriors \\
$\alpha_\lambda,\beta_\lambda$ & shape and rate of the Gamma prior on $\lambda_w,\lambda_v$ \\
$c_0,d_0$ & shape and rate of the Gamma hyperprior on $b_0$ \\
$\theta$ & all model parameters $\{w_0,\{w_j\},\{v_{jf}\}\}$ \\
$S$ & number of retained (post-burn-in) posterior samples \\
$\yhat^{(s)}$ & prediction from posterior sample $s$ \\
\bottomrule
\end{tabular}
\end{table}

We write $\Norm(\mu,\sigma^2)$ for a Gaussian with mean $\mu$ and variance
$\sigma^2$, and $\Gam(a,b)$ for a Gamma distribution with shape $a$ and rate $b$
(so that its mean is $a/b$). All precisions are inverse variances.

% =====================================================================
\subsection{Data representation and feature encoding}
\label{sec:features}

Each observation is a single laboratory LC50 measurement for a
(chemical, species, duration) combination. The target $y_i$ is the
log-transformed concentration in mg/L, centred by subtracting the global mean
$\bar{y}$ of the training set; the mean is added back to recover log-mg/L units
when reporting predictions. Repeated measurements of the same triplet are kept
as separate rows rather than averaged, so that the per-chemical noise model
(Section~\ref{sec:model}) sees the within-triplet spread directly.

The feature vector $\xx_i$ concatenates five one-hot categorical blocks and two
standardised numerical features:
\begin{itemize}
  \item \textbf{chemical} identity (one-hot over $C$ chemicals),
  \item \textbf{species} identity (one-hot),
  \item \textbf{test duration} (one-hot over $\{24,48,72,96\}$\,h),
  \item \textbf{taxonomic family} and \textbf{taxonomic class} of the species (two one-hot blocks),
  \item \textbf{numerical}: log-transformed molecular weight and the RDKit
        octanol--water partition coefficient (\texttt{clogp}), median-imputed
        and standardised to zero mean and unit variance.
\end{itemize}
Unseen categories at prediction time map to an all-zero block (they contribute
nothing to the sums below), so the model can score any (chemical, species,
duration) triplet whose chemical and species were each observed at least once
during training, but not entirely novel chemicals or species. For the dataset
used here $n=70{,}670$, $C=3{,}295$ chemicals, $1{,}267$ species, and
$p=4{,}879$ (namely $3{,}295+1{,}267+4+283+28+2$).

% =====================================================================
\subsection{The factorization-machine predictor}
\label{sec:fm}

The predictive mean for observation $i$ is the second-order factorization
machine
\begin{equation}
  \yhat_i(\theta) \;=\; w_0 \;+\; \sum_{j=1}^{p} w_j\,x_{ij}
  \;+\; \sum_{j=1}^{p}\sum_{j'>j}^{p} \langle \vv_j,\vv_{j'}\rangle\,x_{ij}x_{ij'},
  \label{eq:fm}
\end{equation}
with $\langle \vv_j,\vv_{j'}\rangle=\sum_{f=1}^{k} v_{jf}v_{j'f}$. The pairwise
interactions are not stored explicitly; they are represented in the low-rank
form of Eq.~\eqref{eq:fm}, which lets the model infer the interaction between a
chemical and a species that were never tested together. Using the standard FM
identity, the interaction sum is evaluated in $O(pk)$ time (in practice $O(k
N_{\mathrm{nz}})$ for a sparse $\xx_i$ with $N_{\mathrm{nz}}$ non-zeros) as
\begin{equation}
  \sum_{j<j'} \langle \vv_j,\vv_{j'}\rangle\,x_{ij}x_{ij'}
  \;=\; \tfrac{1}{2}\sum_{f=1}^{k}\!\left[
      \Big(\sum_{j=1}^{p} v_{jf}x_{ij}\Big)^{\!2}
      - \sum_{j=1}^{p} v_{jf}^{2}x_{ij}^{2}
  \right].
  \label{eq:fm-identity}
\end{equation}
The diagonal subtraction in Eq.~\eqref{eq:fm-identity} removes each feature's
self-interaction. A direct consequence, used repeatedly by the sampler, is that
$\yhat_i$ is an \emph{affine} function of each individual parameter when the
others are held fixed (Section~\ref{sec:sampler}).

% =====================================================================
\subsection{Full hierarchical generative model}
\label{sec:model}

The model is a Bayesian FM in which the single global noise precision of the
standard BFM is replaced by a per-chemical precision $\alpha_c$ under a shared
Gamma prior. The complete generative model is:

\paragraph{Likelihood.} Each observation is its FM mean plus Gaussian noise
whose precision is that of the observation's chemical,
\begin{equation}
  y_i \;\mid\; \xx_i,\theta,\alpha_{c_i}
  \;\sim\; \Norm\!\big(\yhat_i(\theta),\ \alpha_{c_i}^{-1}\big),
  \qquad i=1,\dots,n.
  \label{eq:like}
\end{equation}

\paragraph{Priors on the model parameters.} The global bias has a fixed Gaussian
prior; the weights and factors are drawn from shared populations whose mean and
precision are themselves learned:
\begin{align}
  w_0     &\sim \Norm(\mu_0,\ \gamma_0^{-1}), \\
  w_j     &\sim \Norm(\mu_w,\ \lambda_w^{-1}), & j&=1,\dots,p, \label{eq:prior-w}\\
  v_{jf}  &\sim \Norm(\mu_v,\ \lambda_v^{-1}), & j&=1,\dots,p,\ f=1,\dots,k. \label{eq:prior-v}
\end{align}

\paragraph{Hyperpriors on the populations.} There is one population for all
linear weights and one for all latent factors (as in the base BFM of
\citet{rendle2011bayesian}); each has a Gaussian hyperprior on its mean and a
Gamma hyperprior on its precision:
\begin{align}
  \mu_w &\sim \Norm(\mu_0,\ \gamma_0^{-1}), & \lambda_w &\sim \Gam(\alpha_\lambda,\ \beta_\lambda), \\
  \mu_v &\sim \Norm(\mu_0,\ \gamma_0^{-1}), & \lambda_v &\sim \Gam(\alpha_\lambda,\ \beta_\lambda).
\end{align}

\paragraph{Hierarchical noise model.} This is the methodological contribution.
Each chemical has its own noise precision, drawn from a Gamma prior whose rate is
shared across all chemicals and is itself given a Gamma hyperprior:
\begin{align}
  \alpha_c &\sim \Gam(a_0,\ b_0), & c&=1,\dots,C, \label{eq:prior-alpha}\\
  b_0      &\sim \Gam(c_0,\ d_0). \label{eq:prior-b0}
\end{align}
The shape $a_0$ is held fixed (set by empirical Bayes, Section~\ref{sec:eb});
the rate $b_0$ is resampled every sweep. Because $b_0$ is common to all
chemicals, it ties the per-chemical precisions together: a chemical with many
observations has its $\alpha_c$ set mainly by the scatter of its own residuals,
whereas a sparsely measured chemical is pooled towards the population-typical
noise level implied by $\Gam(a_0,b_0)$. This partial pooling is what keeps the
aleatoric estimate stable for data-poor chemicals.

The joint density factorises as
\begin{equation}
\begin{aligned}
  p\big(\theta,\{\alpha_c\},b_0,\mu_w,\lambda_w,\mu_v,\lambda_v \mid \mathbf{y}\big)
  \;\propto\;
  & \Big[\textstyle\prod_{i=1}^{n} \Norm\!\big(y_i \mid \yhat_i(\theta),\alpha_{c_i}^{-1}\big)\Big]
    \Norm(w_0\mid\mu_0,\gamma_0^{-1}) \\
  &\times \Big[\textstyle\prod_{j=1}^{p}\Norm(w_j\mid\mu_w,\lambda_w^{-1})\Big]
          \Big[\textstyle\prod_{j,f}\Norm(v_{jf}\mid\mu_v,\lambda_v^{-1})\Big] \\
  &\times \Norm(\mu_w\mid\mu_0,\gamma_0^{-1})\,\Gam(\lambda_w\mid\alpha_\lambda,\beta_\lambda)\,
          \Norm(\mu_v\mid\mu_0,\gamma_0^{-1})\,\Gam(\lambda_v\mid\alpha_\lambda,\beta_\lambda) \\
  &\times \Big[\textstyle\prod_{c=1}^{C}\Gam(\alpha_c\mid a_0,b_0)\Big]\,\Gam(b_0\mid c_0,d_0).
\end{aligned}
\label{eq:joint}
\end{equation}
This is the model drawn in plate notation in the main text: the fixed constants
$(\mu_0,\gamma_0,\alpha_\lambda,\beta_\lambda,a_0,c_0,d_0)$ set the hyperpriors;
$w_j$ and $v_{jf}$ share $(\mu_w,\lambda_w)$ and $(\mu_v,\lambda_v)$
respectively; the $\alpha_c$ share the rate $b_0$; and each $y_i$ is drawn from a
Gaussian centred at $\yhat_i$ with precision $\alpha_{c_i}$.

% =====================================================================
\subsection{Relationship to the standard Bayesian FM}
\label{sec:diff}

Our model departs from the BFM of \citet{rendle2011bayesian} in two respects.
\begin{enumerate}
  \item \textbf{Per-chemical noise (the substantive change).} The standard BFM
  assumes a single global precision $\alpha$ shared by all observations, with a
  Gamma prior and a Gamma-conjugate posterior. We replace it by $C$
  chemical-specific precisions $\alpha_c$ (Eq.~\ref{eq:prior-alpha}) under a
  shared, learned rate $b_0$ (Eq.~\ref{eq:prior-b0}). Equivalently, the global
  $\alpha$ in every conditional below becomes the observation-specific
  $\alpha_{c_i}$, and one extra conjugate draw for $b_0$ is added per sweep. This
  turns the noise term into a per-chemical, partially pooled quantity, which is
  the aleatoric component carried into the downstream analysis.

  \item \textbf{Independent Gaussian priors on the population means (a minor
  change).} The standard BFM couples each population mean to its precision
  through a Normal--Gamma prior, $\mu \mid \lambda \sim
  \Norm(\mu_0,(\gamma_0\lambda)^{-1})$, so the posterior precision of $\mu$
  scales with $\lambda$. We instead place an independent
  $\Norm(\mu_0,\gamma_0^{-1})$ prior on each mean (Eqs.~\ref{eq:prior-w}--\ref{eq:prior-v}),
  decoupling $\mu$ from $\lambda$; the mean does not then enter the update for
  $\lambda$, and the $\mu$-posterior precision is $\gamma_0+p\lambda$ rather than
  $(\gamma_0+p)\lambda$ (Section~\ref{sec:sampler}). With $p$ in the thousands
  this has negligible effect on inference, as \citet{rendle2011bayesian} also
  note for the hyperprior constants, but we state it for exactness.
\end{enumerate}
As in the reference, we use single-parameter (unblocked) Gibbs sampling, which
avoids the matrix inversions of blocked schemes and keeps each sweep linear in
$k$.

% =====================================================================
\subsection{Gibbs sampler and conditional posteriors}
\label{sec:sampler}

Posterior inference is by Gibbs sampling: each quantity is drawn in turn from its
full conditional given the current values of the others. All conditionals are
available in closed form because (i) the FM mean is affine in each single
parameter, giving Gaussian conditionals with Gaussian priors, and (ii) the
Gamma priors on the precisions are conjugate to the Gaussian likelihood.

Throughout we cache the running residual
\begin{equation}
  e_i \;=\; y_i - \yhat_i(\theta),
  \label{eq:resid}
\end{equation}
and the per-factor cache $q_{if}=\sum_{j=1}^{p} v_{jf}x_{ij}$, both updated by a
rank-one correction after each parameter is resampled. This is what makes a full
sweep $O(k N_{\mathrm{nz}})$, where $N_{\mathrm{nz}}$ is the number of non-zeros
in the design matrix. For a parameter with current value $\theta_\ell$ and
``slope'' $h_{i\ell}=\partial \yhat_i/\partial\theta_\ell$ (constant in
$\theta_\ell$ by multilinearity), the partial residual that adds back its current
contribution is $r_{i}=e_i+h_{i\ell}\theta_\ell$.

\subsubsection{Global bias $w_0$}
Here $h_{i0}=1$ for all $i$, so with prior $\Norm(\mu_0,\gamma_0^{-1})$,
\begin{equation}
  \sigma_{w_0}^2=\Big(\gamma_0+\sum_{i=1}^{n}\alpha_{c_i}\Big)^{-1},
  \qquad
  \mu_{w_0}=\sigma_{w_0}^2\Big(\gamma_0\mu_0+\sum_{i=1}^{n}\alpha_{c_i}\,(e_i+w_0)\Big),
  \qquad
  w_0 \sim \Norm(\mu_{w_0},\sigma_{w_0}^2).
\end{equation}
Then set $e_i \mathrel{+}= (w_0^{\mathrm{old}}-w_0^{\mathrm{new}})$ for all $i$.

\subsubsection{Weight population $(\mu_w,\lambda_w)$}
Treating the $p$ weights as i.i.d.\ $\Norm(\mu_w,\lambda_w^{-1})$ draws,
\begin{align}
  \mu_w &\sim \Norm\!\big(\mu_{\mu_w},\,\sigma_{\mu_w}^2\big),
  & \sigma_{\mu_w}^2 &=\big(\gamma_0+p\,\lambda_w\big)^{-1},
  & \mu_{\mu_w} &=\sigma_{\mu_w}^2\big(\gamma_0\mu_0+\lambda_w\textstyle\sum_{j=1}^{p}w_j\big), \\
  \lambda_w &\sim \Gam\!\Big(\alpha_\lambda+\tfrac{p}{2},\ \
     \beta_\lambda+\tfrac{1}{2}\textstyle\sum_{j=1}^{p}(w_j-\mu_w)^2\Big).
\end{align}

\subsubsection{Linear weights $w_j$}
For feature $j$, $h_{ij}=x_{ij}$, and only observations with $x_{ij}\neq 0$
contribute. With prior $\Norm(\mu_w,\lambda_w^{-1})$ and partial residual
$r_i=e_i+x_{ij}w_j$,
\begin{equation}
  \sigma_{w_j}^2=\Big(\lambda_w+\sum_{i:\,x_{ij}\neq0}\!\alpha_{c_i}\,x_{ij}^2\Big)^{-1},
  \quad
  \mu_{w_j}=\sigma_{w_j}^2\Big(\lambda_w\mu_w+\sum_{i:\,x_{ij}\neq0}\!\alpha_{c_i}\,x_{ij}\,r_i\Big),
  \quad
  w_j\sim\Norm(\mu_{w_j},\sigma_{w_j}^2).
\end{equation}
Then update $e_i \mathrel{+}= x_{ij}\,(w_j^{\mathrm{old}}-w_j^{\mathrm{new}})$.

\subsubsection{Factor population $(\mu_v,\lambda_v)$}
Treating all $pk$ factor entries as i.i.d.\ $\Norm(\mu_v,\lambda_v^{-1})$,
\begin{align}
  \mu_v &\sim \Norm\!\big(\mu_{\mu_v},\,\sigma_{\mu_v}^2\big),
  & \sigma_{\mu_v}^2 &=\big(\gamma_0+p\,k\,\lambda_v\big)^{-1},
  & \mu_{\mu_v} &=\sigma_{\mu_v}^2\big(\gamma_0\mu_0+\lambda_v\textstyle\sum_{j,f}v_{jf}\big), \\
  \lambda_v &\sim \Gam\!\Big(\alpha_\lambda+\tfrac{p k}{2},\ \
     \beta_\lambda+\tfrac{1}{2}\textstyle\sum_{j,f}(v_{jf}-\mu_v)^2\Big).
\end{align}

\subsubsection{Latent factors $v_{jf}$}
By the multilinearity of Eq.~\eqref{eq:fm-identity}, $\yhat_i$ is affine in
$v_{jf}$ with slope
\begin{equation}
  h_{ijf}=x_{ij}\big(q_{if}-x_{ij}v_{jf}\big),
  \qquad q_{if}=\sum_{j'=1}^{p} v_{j'f}x_{ij'},
\end{equation}
which does not depend on $v_{jf}$ (the $x_{ij}v_{jf}$ subtraction removes it).
With prior $\Norm(\mu_v,\lambda_v^{-1})$ and partial residual
$r_i=e_i+h_{ijf}v_{jf}$,
\begin{equation}
  \sigma_{v_{jf}}^2=\Big(\lambda_v+\sum_{i:\,x_{ij}\neq0}\!\alpha_{c_i}\,h_{ijf}^2\Big)^{-1},
  \quad
  \mu_{v_{jf}}=\sigma_{v_{jf}}^2\Big(\lambda_v\mu_v+\sum_{i:\,x_{ij}\neq0}\!\alpha_{c_i}\,h_{ijf}\,r_i\Big),
  \quad
  v_{jf}\sim\Norm(\mu_{v_{jf}},\sigma_{v_{jf}}^2).
\end{equation}
Then update both caches: $e_i \mathrel{+}= h_{ijf}\,(v_{jf}^{\mathrm{old}}-v_{jf}^{\mathrm{new}})$
and $q_{if}\mathrel{+}= x_{ij}\,(v_{jf}^{\mathrm{new}}-v_{jf}^{\mathrm{old}})$.

\subsubsection{Per-chemical noise precisions $\alpha_c$}
Let $N_c=|\{i:c_i=c\}|$ and $\mathrm{SSE}_c=\sum_{i:\,c_i=c} e_i^2$ be the
observation count and residual sum of squares of chemical $c$. Conjugacy of the
Gamma prior (Eq.~\ref{eq:prior-alpha}) to the Gaussian likelihood
(Eq.~\ref{eq:like}) gives
\begin{equation}
  \alpha_c \;\sim\; \Gam\!\Big(a_0+\tfrac{N_c}{2},\ \ b_0+\tfrac{1}{2}\,\mathrm{SSE}_c\Big),
  \qquad c=1,\dots,C.
  \label{eq:post-alpha}
\end{equation}
All $C$ precisions are drawn in one vectorised step using per-group sums of
$e_i^2$.

\subsubsection{Shared noise rate $b_0$}
The rate $b_0$ is conjugate to the $C$ Gamma-distributed precisions with common
shape $a_0$: combining $\prod_c \Gam(\alpha_c\mid a_0,b_0)\propto
b_0^{C a_0}e^{-b_0\sum_c\alpha_c}$ with the prior $\Gam(c_0,d_0)$ yields
\begin{equation}
  b_0 \;\sim\; \Gam\!\Big(c_0+C\,a_0,\ \ d_0+\textstyle\sum_{c=1}^{C}\alpha_c\Big).
  \label{eq:post-b0}
\end{equation}

\subsubsection{One sweep and sample collection}
A single Gibbs sweep applies, in order: $w_0$; $(\mu_w,\lambda_w)$ then every
$w_j$; $(\mu_v,\lambda_v)$ then every $v_{jf}$; the $\alpha_c$
(Eq.~\ref{eq:post-alpha}); and $b_0$ (Eq.~\ref{eq:post-b0}). The chain is run
for $n_{\mathrm{iter}}$ sweeps from an arbitrary initialisation
($w_0=0$, $w=\mathbf{0}$, $v_{jf}\sim\Norm(0,0.1^2)$, $\alpha_c=1$); the first
$n_{\mathrm{burn}}$ sweeps are discarded as burn-in and the remaining
$S=n_{\mathrm{iter}}-n_{\mathrm{burn}}$ sweeps are retained. Each retained sample
$s$ stores $\{w_0^{(s)},w^{(s)},v^{(s)},\{\alpha_c^{(s)}\},b_0^{(s)}\}$ and
corresponds to one complete, internally consistent parameterisation of the
model.

% =====================================================================
\subsection{Empirical-Bayes initialisation of the prior shape $a_0$}
\label{sec:eb}

The shared shape $a_0$ controls how strongly sparsely tested chemicals are
pooled towards the population noise level. Rather than fix it arbitrarily, we set
it by a method-of-moments empirical-Bayes pre-pass so that the prior matches the
dispersion of \emph{well-estimated} precisions:
\begin{enumerate}
  \item Run a short Gibbs pre-pass on the full training set with a non-informative,
  fixed noise prior ($a_0=1$, $b_0=1$, $b_0$ not learned) and collect the
  posterior draws $\{\alpha_c^{(s)}\}$.
  \item Restrict to data-rich chemicals ($N_c\ge N_{\min}$, with $N_{\min}=50$),
  whose posterior precisions are essentially data-driven, and form their
  posterior means $\hat{\alpha}_c=\tfrac{1}{S}\sum_s\alpha_c^{(s)}$.
  \item Fit a Gamma to $\{\hat{\alpha}_c\}$ by matching mean $m$ and variance
  $\nu$, and take the shape
  \begin{equation}
    a_0 \;=\; m^2/\nu.
  \end{equation}
  The corresponding rate is discarded, because the main run learns $b_0$ from its
  hyperprior (Eq.~\ref{eq:post-b0}).
\end{enumerate}
For the run reported in the paper this gives $a_0\approx1.75$; the main-run
posterior for the shared rate concentrates near $b_0\approx0.08$, so
$\Gam(a_0,b_0)$ is a weakly informative population prior on the per-chemical
precisions. The same EB $a_0$ is used for the cross-validation runs and for the
full-data model used to generate SSDs.

% =====================================================================
\subsection{Hyperparameter settings}
\label{sec:hyper}

Table~\ref{tab:hyper} lists the fixed constants and run settings. The mean- and
weight/factor-precision hyperpriors are the same trivial constants used by
\citet{rendle2011bayesian}; with $p$ in the thousands their influence on the
posterior is negligible. The noise-model hyperprior $\Gam(c_0,d_0)$ on $b_0$ and
the EB shape $a_0$ are specific to this work.

\begin{table}[htbp]
\centering
\caption{Hyperparameter and run settings.}
\label{tab:hyper}
\small
\begin{tabular}{lll}
\toprule
\textbf{Quantity} & \textbf{Symbol} & \textbf{Value} \\
\midrule
Factorization rank            & $k$                       & $32$ \\
Mean-hyperprior mean          & $\mu_0$                   & $0$ \\
Mean-hyperprior precision     & $\gamma_0$                & $1$ \\
Weight/factor precision prior & $\alpha_\lambda,\beta_\lambda$ & $1,\ 1$ \\
Noise-rate hyperprior         & $c_0,\ d_0$               & $1,\ 1$ \\
Noise-prior shape (empirical Bayes) & $a_0$               & $\approx 1.75$ (fixed after EB) \\
Noise-prior rate              & $b_0$                     & initialised at $1$, learned \\
Main run                      & $n_{\mathrm{iter}},\ n_{\mathrm{burn}}$ & $2000,\ 100$ \\
Retained samples              & $S$                       & $1900$ \\
EB pre-pass                   & $n_{\mathrm{iter}},\ n_{\mathrm{burn}}$ & $100,\ 50$ \\
EB data-rich threshold        & $N_{\min}$                & $50$ \\
Cross-validation              & folds                     & $5$ (grouped by triplet) \\
\bottomrule
\end{tabular}
\end{table}

% =====================================================================
\subsection{Posterior-predictive quantities}
\label{sec:predictive}

For a query $\xx^\ast$ belonging to chemical $c^\ast$, each retained sample gives
a deterministic FM prediction $\yhat^{(s)}(\xx^\ast)$ via Eq.~\eqref{eq:fm}. The
point prediction is the sample mean
\begin{equation}
  \yhat^\ast=\frac{1}{S}\sum_{s=1}^{S}\yhat^{(s)}(\xx^\ast).
\end{equation}
The predictive variance decomposes by the law of total variance,
$\Var(y^\ast\mid\mathcal{D})=\Var_\theta\!\big(\E[y^\ast\mid\theta]\big)
+\E_\theta\!\big[\Var(y^\ast\mid\theta)\big]$, whose two Monte-Carlo estimates are
the epistemic and aleatoric components:
\begin{align}
  \Var_{\mathrm{epi}}(\xx^\ast) &= \frac{1}{S}\sum_{s=1}^{S}\big(\yhat^{(s)}(\xx^\ast)-\yhat^\ast\big)^2,
  \label{eq:epi}\\
  \Var_{\mathrm{ale}}(c^\ast) &= \frac{1}{S}\sum_{s=1}^{S}\frac{1}{\alpha_{c^\ast}^{(s)}},
  \label{eq:ale}\\
  \Var_{\mathrm{tot}}(\xx^\ast,c^\ast) &= \Var_{\mathrm{epi}}(\xx^\ast)+\Var_{\mathrm{ale}}(c^\ast).
  \label{eq:tot}
\end{align}
The epistemic term is the spread of the mean across plausible fits and shrinks as
data constrain the parameters; the aleatoric term is the posterior-mean noise
variance of the chemical and does not vanish with more data. The aleatoric term
depends only on the chemical, so it is constant across species and durations for
a given chemical.

\paragraph{Calibration.} To assess whether $\Var_{\mathrm{tot}}$ is well
calibrated we form posterior-predictive draws that combine both components. For
each held-out observation and each sample $s$ we draw
\begin{equation}
  \tilde{y}^{(s)} \sim \Norm\!\big(\yhat^{(s)}(\xx^\ast),\ 1/\alpha_{c^\ast}^{(s)}\big),
\end{equation}
giving an $S$-point predictive sample. For a nominal central level $\rho$ we take
the empirical $[(1-\rho)/2,\ (1+\rho)/2]$ quantiles of these draws as the
credible interval, and report empirical coverage as the fraction of held-out
observations whose true value lies inside. Coverage is evaluated on the
out-of-fold observations (Section~\ref{sec:cv}), so predictions are scored by
models that did not train on them.

% =====================================================================
\subsection{SSD ensembles and hazardous concentrations}
\label{sec:ssd}

The Bayesian structure lets uncertainty propagate coherently into Species
Sensitivity Distributions. For a target chemical $c^\ast$ at a fixed duration,
and each retained sample $s$, we evaluate Eq.~\eqref{eq:fm} at $\theta^{(s)}$ for
all species $i=1,\dots,M$ and sort the resulting predictions ascending to obtain
one SSD curve $\big(\yhat^{(s)}_{(1)}\le\cdots\le\yhat^{(s)}_{(M)}\big)$ plotted
against the plotting positions $m/(M+1)$. Because every curve comes from a single
coherent parameter draw, the correlated epistemic structure across species is
preserved: the whole curve shifts and rescales as the shared parameters vary,
rather than species scattering independently. The ensemble of $S$ curves is
summarised pointwise by its median and central $95\%$ band.

A hazardous concentration $\mathrm{HC}x$ (the concentration protective of
$100-x\%$ of species) is read off each sampled curve at rank
$\lfloor x\,M/100\rfloor$, giving a posterior sample of $\mathrm{HC}x$ values;
its median and $2.5/97.5$ percentiles are the reported estimate and $95\%$
credible interval. Optionally, aleatoric noise can be injected per species by
adding a draw $\Norm\!\big(0,1/\alpha_{c^\ast}^{(s)}\big)$ to each prediction
before sorting; the SSD figures and HC intervals in the main text use the
epistemic ensemble (no added aleatoric noise) unless stated otherwise.

% =====================================================================
\subsection{Cross-validation and out-of-fold uncertainty}
\label{sec:cv}

Predictive accuracy and calibration are evaluated with grouped $5$-fold
cross-validation. Folds are formed by \emph{grouped} splitting on the
(chemical, species, duration) triplet identifier, so that all replicate
measurements of a held-out triplet are removed from training together and no
replicate of a scored triplet is ever seen during training. Each fold refits the
full model (including the hierarchical noise model, with the EB shape $a_0$ fixed
and $b_0$ learned) on its training partition and predicts the held-out partition;
concatenating the held-out predictions over folds yields the out-of-fold arrays
used for the RMSE, the uncertainty-versus-data analyses
(Eqs.~\ref{eq:epi}--\ref{eq:ale}), and the calibration curve. The full-data model
used to generate the SSDs and the all-chemical HC comparison is fit once on all
$n$ observations with the same settings.

Two practical points. First, the epistemic estimate of Eq.~\eqref{eq:epi} is only
meaningful for chemicals that appear in the training partition of a fold;
chemicals with fewer total observations than the number of folds can have folds
in which none of their rows are present, leaving their chemical-specific
parameters at the prior and deflating their apparent epistemic variance. Such
cold-start chemicals are excluded from the per-chemical
uncertainty-versus-data analyses in the main text. Second, the aleatoric
out-of-fold value for an observation is the sample mean of
$1/\alpha_{c}^{(s)}$ over that fold's retained samples, so it reflects the noise
precision learned from the fold's training data only.

% =====================================================================
\subsection{Computational complexity}
\label{sec:complexity}

Every conditional in Section~\ref{sec:sampler} is a single Gaussian or Gamma
draw, and the residual and factor caches make each parameter update a rank-one
correction touching only the non-zeros of the relevant feature column. One full
sweep therefore costs $O(k\,N_{\mathrm{nz}})$, linear in the factorization rank
and in the number of non-zero design-matrix entries, with an additional
$O(C)$ for the noise block (Eqs.~\ref{eq:post-alpha}--\ref{eq:post-b0}). This
matches the complexity of the unblocked BFM sampler of
\citet{rendle2011bayesian}; the per-chemical noise model adds only the vectorised
$\alpha_c$ draw and one scalar $b_0$ draw per sweep, so the hierarchical
extension does not change the asymptotic cost.

.
% Relies on the preamble of main.tex (amsmath, booktabs, etc. are already
% loaded there; amssymb is added there for \Real below). Uses the shared
% references.bib already cited by the rest of the paper.
% =====================================================================

% ---- convenience macros ----
\newcommand{\yhat}{\hat{y}}
\newcommand{\Norm}{\mathcal{N}}
\newcommand{\Gam}{\mathrm{Gamma}}
\newcommand{\Real}{\mathbb{R}}
\newcommand{\E}{\mathbb{E}}
\newcommand{\Var}{\mathrm{Var}}
\newcommand{\vv}{\mathbf{v}}
\newcommand{\xx}{\mathbf{x}}

\section{Supporting Information}
\label{sec:si}

This Supporting Information gives the full mathematical specification of the
Hierarchical Bayesian Factorization Machine (BFM) summarised in the main text.
It is written for a statistical or machine-learning audience and documents the
model exactly as implemented: the predictor, the complete hierarchical
generative model, the Gibbs sampler and every conditional posterior, the
empirical-Bayes initialisation, the hyperparameter settings, and the
posterior-predictive quantities used for calibration, Species Sensitivity
Distributions (SSDs), and hazardous concentrations. The predictive core is the
factorization machine (FM) of \citet{rendle2010fm} cast in the Bayesian form of
\citet{rendle2011bayesian}; Section~\ref{sec:diff} states precisely how our
model differs from theirs. Symbols follow the main text and are collected in
Table~\ref{tab:notation}.

% =====================================================================
\subsection{Notation}
\label{sec:notation}

\begin{table}[htbp]
\centering
\caption{Notation used throughout the Supporting Information.}
\label{tab:notation}
\small
\begin{tabular}{ll}
\toprule
\textbf{Symbol} & \textbf{Meaning} \\
\midrule
$n$ & number of observations (rows) \\
$p$ & number of features (design-matrix columns) \\
$k$ & number of latent factors per feature (factorization rank) \\
$C$ & number of chemicals (noise groups) \\
$\xx_i \in \Real^{p}$ & feature vector of observation $i$ \\
$x_{ij}$ & $j$-th entry of $\xx_i$ \\
$y_i$ & centred log-transformed LC50 of observation $i$ \\
$c_i \in \{1,\dots,C\}$ & chemical (group) index of observation $i$ \\
$w_0$ & global bias \\
$w_j$ & linear weight of feature $j$ \\
$v_{jf}$ & $f$-th latent factor of feature $j$; $\vv_j=(v_{j1},\dots,v_{jk})$ \\
$\alpha_c$ & noise precision of chemical $c$ (aleatoric term) \\
$\mu_w,\lambda_w$ & mean and precision of the shared prior on the weights $w_j$ \\
$\mu_v,\lambda_v$ & mean and precision of the shared prior on the factors $v_{jf}$ \\
$a_0,b_0$ & shape and rate of the shared prior on the $\alpha_c$ \\
$\mu_0,\gamma_0$ & mean and precision of the Gaussian mean-hyperpriors \\
$\alpha_\lambda,\beta_\lambda$ & shape and rate of the Gamma prior on $\lambda_w,\lambda_v$ \\
$c_0,d_0$ & shape and rate of the Gamma hyperprior on $b_0$ \\
$\theta$ & all model parameters $\{w_0,\{w_j\},\{v_{jf}\}\}$ \\
$S$ & number of retained (post-burn-in) posterior samples \\
$\yhat^{(s)}$ & prediction from posterior sample $s$ \\
\bottomrule
\end{tabular}
\end{table}

We write $\Norm(\mu,\sigma^2)$ for a Gaussian with mean $\mu$ and variance
$\sigma^2$, and $\Gam(a,b)$ for a Gamma distribution with shape $a$ and rate $b$
(so that its mean is $a/b$). All precisions are inverse variances.

% =====================================================================
\subsection{Data representation and feature encoding}
\label{sec:features}

Each observation is a single laboratory LC50 measurement for a
(chemical, species, duration) combination. The target $y_i$ is the
log-transformed concentration in mg/L, centred by subtracting the global mean
$\bar{y}$ of the training set; the mean is added back to recover log-mg/L units
when reporting predictions. Repeated measurements of the same triplet are kept
as separate rows rather than averaged, so that the per-chemical noise model
(Section~\ref{sec:model}) sees the within-triplet spread directly.

The feature vector $\xx_i$ concatenates five one-hot categorical blocks and two
standardised numerical features:
\begin{itemize}
  \item \textbf{chemical} identity (one-hot over $C$ chemicals),
  \item \textbf{species} identity (one-hot),
  \item \textbf{test duration} (one-hot over $\{24,48,72,96\}$\,h),
  \item \textbf{taxonomic family} and \textbf{taxonomic class} of the species (two one-hot blocks),
  \item \textbf{numerical}: log-transformed molecular weight and the RDKit
        octanol--water partition coefficient (\texttt{clogp}), median-imputed
        and standardised to zero mean and unit variance.
\end{itemize}
Unseen categories at prediction time map to an all-zero block (they contribute
nothing to the sums below), so the model can score any (chemical, species,
duration) triplet whose chemical and species were each observed at least once
during training, but not entirely novel chemicals or species. For the dataset
used here $n=70{,}670$, $C=3{,}295$ chemicals, $1{,}267$ species, and
$p=4{,}879$ (namely $3{,}295+1{,}267+4+283+28+2$).

% =====================================================================
\subsection{The factorization-machine predictor}
\label{sec:fm}

The predictive mean for observation $i$ is the second-order factorization
machine
\begin{equation}
  \yhat_i(\theta) \;=\; w_0 \;+\; \sum_{j=1}^{p} w_j\,x_{ij}
  \;+\; \sum_{j=1}^{p}\sum_{j'>j}^{p} \langle \vv_j,\vv_{j'}\rangle\,x_{ij}x_{ij'},
  \label{eq:fm}
\end{equation}
with $\langle \vv_j,\vv_{j'}\rangle=\sum_{f=1}^{k} v_{jf}v_{j'f}$. The pairwise
interactions are not stored explicitly; they are represented in the low-rank
form of Eq.~\eqref{eq:fm}, which lets the model infer the interaction between a
chemical and a species that were never tested together. Using the standard FM
identity, the interaction sum is evaluated in $O(pk)$ time (in practice $O(k
N_{\mathrm{nz}})$ for a sparse $\xx_i$ with $N_{\mathrm{nz}}$ non-zeros) as
\begin{equation}
  \sum_{j<j'} \langle \vv_j,\vv_{j'}\rangle\,x_{ij}x_{ij'}
  \;=\; \tfrac{1}{2}\sum_{f=1}^{k}\!\left[
      \Big(\sum_{j=1}^{p} v_{jf}x_{ij}\Big)^{\!2}
      - \sum_{j=1}^{p} v_{jf}^{2}x_{ij}^{2}
  \right].
  \label{eq:fm-identity}
\end{equation}
The diagonal subtraction in Eq.~\eqref{eq:fm-identity} removes each feature's
self-interaction. A direct consequence, used repeatedly by the sampler, is that
$\yhat_i$ is an \emph{affine} function of each individual parameter when the
others are held fixed (Section~\ref{sec:sampler}).

% =====================================================================
\subsection{Full hierarchical generative model}
\label{sec:model}

The model is a Bayesian FM in which the single global noise precision of the
standard BFM is replaced by a per-chemical precision $\alpha_c$ under a shared
Gamma prior. The complete generative model is:

\paragraph{Likelihood.} Each observation is its FM mean plus Gaussian noise
whose precision is that of the observation's chemical,
\begin{equation}
  y_i \;\mid\; \xx_i,\theta,\alpha_{c_i}
  \;\sim\; \Norm\!\big(\yhat_i(\theta),\ \alpha_{c_i}^{-1}\big),
  \qquad i=1,\dots,n.
  \label{eq:like}
\end{equation}

\paragraph{Priors on the model parameters.} The global bias has a fixed Gaussian
prior; the weights and factors are drawn from shared populations whose mean and
precision are themselves learned:
\begin{align}
  w_0     &\sim \Norm(\mu_0,\ \gamma_0^{-1}), \\
  w_j     &\sim \Norm(\mu_w,\ \lambda_w^{-1}), & j&=1,\dots,p, \label{eq:prior-w}\\
  v_{jf}  &\sim \Norm(\mu_v,\ \lambda_v^{-1}), & j&=1,\dots,p,\ f=1,\dots,k. \label{eq:prior-v}
\end{align}

\paragraph{Hyperpriors on the populations.} There is one population for all
linear weights and one for all latent factors (as in the base BFM of
\citet{rendle2011bayesian}); each has a Gaussian hyperprior on its mean and a
Gamma hyperprior on its precision:
\begin{align}
  \mu_w &\sim \Norm(\mu_0,\ \gamma_0^{-1}), & \lambda_w &\sim \Gam(\alpha_\lambda,\ \beta_\lambda), \\
  \mu_v &\sim \Norm(\mu_0,\ \gamma_0^{-1}), & \lambda_v &\sim \Gam(\alpha_\lambda,\ \beta_\lambda).
\end{align}

\paragraph{Hierarchical noise model.} This is the methodological contribution.
Each chemical has its own noise precision, drawn from a Gamma prior whose rate is
shared across all chemicals and is itself given a Gamma hyperprior:
\begin{align}
  \alpha_c &\sim \Gam(a_0,\ b_0), & c&=1,\dots,C, \label{eq:prior-alpha}\\
  b_0      &\sim \Gam(c_0,\ d_0). \label{eq:prior-b0}
\end{align}
The shape $a_0$ is held fixed (set by empirical Bayes, Section~\ref{sec:eb});
the rate $b_0$ is resampled every sweep. Because $b_0$ is common to all
chemicals, it ties the per-chemical precisions together: a chemical with many
observations has its $\alpha_c$ set mainly by the scatter of its own residuals,
whereas a sparsely measured chemical is pooled towards the population-typical
noise level implied by $\Gam(a_0,b_0)$. This partial pooling is what keeps the
aleatoric estimate stable for data-poor chemicals.

The joint density factorises as
\begin{equation}
\begin{aligned}
  p\big(\theta,\{\alpha_c\},b_0,\mu_w,\lambda_w,\mu_v,\lambda_v \mid \mathbf{y}\big)
  \;\propto\;
  & \Big[\textstyle\prod_{i=1}^{n} \Norm\!\big(y_i \mid \yhat_i(\theta),\alpha_{c_i}^{-1}\big)\Big]
    \Norm(w_0\mid\mu_0,\gamma_0^{-1}) \\
  &\times \Big[\textstyle\prod_{j=1}^{p}\Norm(w_j\mid\mu_w,\lambda_w^{-1})\Big]
          \Big[\textstyle\prod_{j,f}\Norm(v_{jf}\mid\mu_v,\lambda_v^{-1})\Big] \\
  &\times \Norm(\mu_w\mid\mu_0,\gamma_0^{-1})\,\Gam(\lambda_w\mid\alpha_\lambda,\beta_\lambda)\,
          \Norm(\mu_v\mid\mu_0,\gamma_0^{-1})\,\Gam(\lambda_v\mid\alpha_\lambda,\beta_\lambda) \\
  &\times \Big[\textstyle\prod_{c=1}^{C}\Gam(\alpha_c\mid a_0,b_0)\Big]\,\Gam(b_0\mid c_0,d_0).
\end{aligned}
\label{eq:joint}
\end{equation}
This is the model drawn in plate notation in the main text: the fixed constants
$(\mu_0,\gamma_0,\alpha_\lambda,\beta_\lambda,a_0,c_0,d_0)$ set the hyperpriors;
$w_j$ and $v_{jf}$ share $(\mu_w,\lambda_w)$ and $(\mu_v,\lambda_v)$
respectively; the $\alpha_c$ share the rate $b_0$; and each $y_i$ is drawn from a
Gaussian centred at $\yhat_i$ with precision $\alpha_{c_i}$.

% =====================================================================
\subsection{Relationship to the standard Bayesian FM}
\label{sec:diff}

Our model departs from the BFM of \citet{rendle2011bayesian} in two respects.
\begin{enumerate}
  \item \textbf{Per-chemical noise (the substantive change).} The standard BFM
  assumes a single global precision $\alpha$ shared by all observations, with a
  Gamma prior and a Gamma-conjugate posterior. We replace it by $C$
  chemical-specific precisions $\alpha_c$ (Eq.~\ref{eq:prior-alpha}) under a
  shared, learned rate $b_0$ (Eq.~\ref{eq:prior-b0}). Equivalently, the global
  $\alpha$ in every conditional below becomes the observation-specific
  $\alpha_{c_i}$, and one extra conjugate draw for $b_0$ is added per sweep. This
  turns the noise term into a per-chemical, partially pooled quantity, which is
  the aleatoric component carried into the downstream analysis.

  \item \textbf{Independent Gaussian priors on the population means (a minor
  change).} The standard BFM couples each population mean to its precision
  through a Normal--Gamma prior, $\mu \mid \lambda \sim
  \Norm(\mu_0,(\gamma_0\lambda)^{-1})$, so the posterior precision of $\mu$
  scales with $\lambda$. We instead place an independent
  $\Norm(\mu_0,\gamma_0^{-1})$ prior on each mean (Eqs.~\ref{eq:prior-w}--\ref{eq:prior-v}),
  decoupling $\mu$ from $\lambda$; the mean does not then enter the update for
  $\lambda$, and the $\mu$-posterior precision is $\gamma_0+p\lambda$ rather than
  $(\gamma_0+p)\lambda$ (Section~\ref{sec:sampler}). With $p$ in the thousands
  this has negligible effect on inference, as \citet{rendle2011bayesian} also
  note for the hyperprior constants, but we state it for exactness.
\end{enumerate}
As in the reference, we use single-parameter (unblocked) Gibbs sampling, which
avoids the matrix inversions of blocked schemes and keeps each sweep linear in
$k$.

% =====================================================================
\subsection{Gibbs sampler and conditional posteriors}
\label{sec:sampler}

Posterior inference is by Gibbs sampling: each quantity is drawn in turn from its
full conditional given the current values of the others. All conditionals are
available in closed form because (i) the FM mean is affine in each single
parameter, giving Gaussian conditionals with Gaussian priors, and (ii) the
Gamma priors on the precisions are conjugate to the Gaussian likelihood.

Throughout we cache the running residual
\begin{equation}
  e_i \;=\; y_i - \yhat_i(\theta),
  \label{eq:resid}
\end{equation}
and the per-factor cache $q_{if}=\sum_{j=1}^{p} v_{jf}x_{ij}$, both updated by a
rank-one correction after each parameter is resampled. This is what makes a full
sweep $O(k N_{\mathrm{nz}})$, where $N_{\mathrm{nz}}$ is the number of non-zeros
in the design matrix. For a parameter with current value $\theta_\ell$ and
``slope'' $h_{i\ell}=\partial \yhat_i/\partial\theta_\ell$ (constant in
$\theta_\ell$ by multilinearity), the partial residual that adds back its current
contribution is $r_{i}=e_i+h_{i\ell}\theta_\ell$.

\subsubsection{Global bias $w_0$}
Here $h_{i0}=1$ for all $i$, so with prior $\Norm(\mu_0,\gamma_0^{-1})$,
\begin{equation}
  \sigma_{w_0}^2=\Big(\gamma_0+\sum_{i=1}^{n}\alpha_{c_i}\Big)^{-1},
  \qquad
  \mu_{w_0}=\sigma_{w_0}^2\Big(\gamma_0\mu_0+\sum_{i=1}^{n}\alpha_{c_i}\,(e_i+w_0)\Big),
  \qquad
  w_0 \sim \Norm(\mu_{w_0},\sigma_{w_0}^2).
\end{equation}
Then set $e_i \mathrel{+}= (w_0^{\mathrm{old}}-w_0^{\mathrm{new}})$ for all $i$.

\subsubsection{Weight population $(\mu_w,\lambda_w)$}
Treating the $p$ weights as i.i.d.\ $\Norm(\mu_w,\lambda_w^{-1})$ draws,
\begin{align}
  \mu_w &\sim \Norm\!\big(\mu_{\mu_w},\,\sigma_{\mu_w}^2\big),
  & \sigma_{\mu_w}^2 &=\big(\gamma_0+p\,\lambda_w\big)^{-1},
  & \mu_{\mu_w} &=\sigma_{\mu_w}^2\big(\gamma_0\mu_0+\lambda_w\textstyle\sum_{j=1}^{p}w_j\big), \\
  \lambda_w &\sim \Gam\!\Big(\alpha_\lambda+\tfrac{p}{2},\ \
     \beta_\lambda+\tfrac{1}{2}\textstyle\sum_{j=1}^{p}(w_j-\mu_w)^2\Big).
\end{align}

\subsubsection{Linear weights $w_j$}
For feature $j$, $h_{ij}=x_{ij}$, and only observations with $x_{ij}\neq 0$
contribute. With prior $\Norm(\mu_w,\lambda_w^{-1})$ and partial residual
$r_i=e_i+x_{ij}w_j$,
\begin{equation}
  \sigma_{w_j}^2=\Big(\lambda_w+\sum_{i:\,x_{ij}\neq0}\!\alpha_{c_i}\,x_{ij}^2\Big)^{-1},
  \quad
  \mu_{w_j}=\sigma_{w_j}^2\Big(\lambda_w\mu_w+\sum_{i:\,x_{ij}\neq0}\!\alpha_{c_i}\,x_{ij}\,r_i\Big),
  \quad
  w_j\sim\Norm(\mu_{w_j},\sigma_{w_j}^2).
\end{equation}
Then update $e_i \mathrel{+}= x_{ij}\,(w_j^{\mathrm{old}}-w_j^{\mathrm{new}})$.

\subsubsection{Factor population $(\mu_v,\lambda_v)$}
Treating all $pk$ factor entries as i.i.d.\ $\Norm(\mu_v,\lambda_v^{-1})$,
\begin{align}
  \mu_v &\sim \Norm\!\big(\mu_{\mu_v},\,\sigma_{\mu_v}^2\big),
  & \sigma_{\mu_v}^2 &=\big(\gamma_0+p\,k\,\lambda_v\big)^{-1},
  & \mu_{\mu_v} &=\sigma_{\mu_v}^2\big(\gamma_0\mu_0+\lambda_v\textstyle\sum_{j,f}v_{jf}\big), \\
  \lambda_v &\sim \Gam\!\Big(\alpha_\lambda+\tfrac{p k}{2},\ \
     \beta_\lambda+\tfrac{1}{2}\textstyle\sum_{j,f}(v_{jf}-\mu_v)^2\Big).
\end{align}

\subsubsection{Latent factors $v_{jf}$}
By the multilinearity of Eq.~\eqref{eq:fm-identity}, $\yhat_i$ is affine in
$v_{jf}$ with slope
\begin{equation}
  h_{ijf}=x_{ij}\big(q_{if}-x_{ij}v_{jf}\big),
  \qquad q_{if}=\sum_{j'=1}^{p} v_{j'f}x_{ij'},
\end{equation}
which does not depend on $v_{jf}$ (the $x_{ij}v_{jf}$ subtraction removes it).
With prior $\Norm(\mu_v,\lambda_v^{-1})$ and partial residual
$r_i=e_i+h_{ijf}v_{jf}$,
\begin{equation}
  \sigma_{v_{jf}}^2=\Big(\lambda_v+\sum_{i:\,x_{ij}\neq0}\!\alpha_{c_i}\,h_{ijf}^2\Big)^{-1},
  \quad
  \mu_{v_{jf}}=\sigma_{v_{jf}}^2\Big(\lambda_v\mu_v+\sum_{i:\,x_{ij}\neq0}\!\alpha_{c_i}\,h_{ijf}\,r_i\Big),
  \quad
  v_{jf}\sim\Norm(\mu_{v_{jf}},\sigma_{v_{jf}}^2).
\end{equation}
Then update both caches: $e_i \mathrel{+}= h_{ijf}\,(v_{jf}^{\mathrm{old}}-v_{jf}^{\mathrm{new}})$
and $q_{if}\mathrel{+}= x_{ij}\,(v_{jf}^{\mathrm{new}}-v_{jf}^{\mathrm{old}})$.

\subsubsection{Per-chemical noise precisions $\alpha_c$}
Let $N_c=|\{i:c_i=c\}|$ and $\mathrm{SSE}_c=\sum_{i:\,c_i=c} e_i^2$ be the
observation count and residual sum of squares of chemical $c$. Conjugacy of the
Gamma prior (Eq.~\ref{eq:prior-alpha}) to the Gaussian likelihood
(Eq.~\ref{eq:like}) gives
\begin{equation}
  \alpha_c \;\sim\; \Gam\!\Big(a_0+\tfrac{N_c}{2},\ \ b_0+\tfrac{1}{2}\,\mathrm{SSE}_c\Big),
  \qquad c=1,\dots,C.
  \label{eq:post-alpha}
\end{equation}
All $C$ precisions are drawn in one vectorised step using per-group sums of
$e_i^2$.

\subsubsection{Shared noise rate $b_0$}
The rate $b_0$ is conjugate to the $C$ Gamma-distributed precisions with common
shape $a_0$: combining $\prod_c \Gam(\alpha_c\mid a_0,b_0)\propto
b_0^{C a_0}e^{-b_0\sum_c\alpha_c}$ with the prior $\Gam(c_0,d_0)$ yields
\begin{equation}
  b_0 \;\sim\; \Gam\!\Big(c_0+C\,a_0,\ \ d_0+\textstyle\sum_{c=1}^{C}\alpha_c\Big).
  \label{eq:post-b0}
\end{equation}

\subsubsection{One sweep and sample collection}
A single Gibbs sweep applies, in order: $w_0$; $(\mu_w,\lambda_w)$ then every
$w_j$; $(\mu_v,\lambda_v)$ then every $v_{jf}$; the $\alpha_c$
(Eq.~\ref{eq:post-alpha}); and $b_0$ (Eq.~\ref{eq:post-b0}). The chain is run
for $n_{\mathrm{iter}}$ sweeps from an arbitrary initialisation
($w_0=0$, $w=\mathbf{0}$, $v_{jf}\sim\Norm(0,0.1^2)$, $\alpha_c=1$); the first
$n_{\mathrm{burn}}$ sweeps are discarded as burn-in and the remaining
$S=n_{\mathrm{iter}}-n_{\mathrm{burn}}$ sweeps are retained. Each retained sample
$s$ stores $\{w_0^{(s)},w^{(s)},v^{(s)},\{\alpha_c^{(s)}\},b_0^{(s)}\}$ and
corresponds to one complete, internally consistent parameterisation of the
model.

% =====================================================================
\subsection{Empirical-Bayes initialisation of the prior shape $a_0$}
\label{sec:eb}

The shared shape $a_0$ controls how strongly sparsely tested chemicals are
pooled towards the population noise level. Rather than fix it arbitrarily, we set
it by a method-of-moments empirical-Bayes pre-pass so that the prior matches the
dispersion of \emph{well-estimated} precisions:
\begin{enumerate}
  \item Run a short Gibbs pre-pass on the full training set with a non-informative,
  fixed noise prior ($a_0=1$, $b_0=1$, $b_0$ not learned) and collect the
  posterior draws $\{\alpha_c^{(s)}\}$.
  \item Restrict to data-rich chemicals ($N_c\ge N_{\min}$, with $N_{\min}=50$),
  whose posterior precisions are essentially data-driven, and form their
  posterior means $\hat{\alpha}_c=\tfrac{1}{S}\sum_s\alpha_c^{(s)}$.
  \item Fit a Gamma to $\{\hat{\alpha}_c\}$ by matching mean $m$ and variance
  $\nu$, and take the shape
  \begin{equation}
    a_0 \;=\; m^2/\nu.
  \end{equation}
  The corresponding rate is discarded, because the main run learns $b_0$ from its
  hyperprior (Eq.~\ref{eq:post-b0}).
\end{enumerate}
For the run reported in the paper this gives $a_0\approx1.75$; the main-run
posterior for the shared rate concentrates near $b_0\approx0.08$, so
$\Gam(a_0,b_0)$ is a weakly informative population prior on the per-chemical
precisions. The same EB $a_0$ is used for the cross-validation runs and for the
full-data model used to generate SSDs.

% =====================================================================
\subsection{Hyperparameter settings}
\label{sec:hyper}

Table~\ref{tab:hyper} lists the fixed constants and run settings. The mean- and
weight/factor-precision hyperpriors are the same trivial constants used by
\citet{rendle2011bayesian}; with $p$ in the thousands their influence on the
posterior is negligible. The noise-model hyperprior $\Gam(c_0,d_0)$ on $b_0$ and
the EB shape $a_0$ are specific to this work.

\begin{table}[htbp]
\centering
\caption{Hyperparameter and run settings.}
\label{tab:hyper}
\small
\begin{tabular}{lll}
\toprule
\textbf{Quantity} & \textbf{Symbol} & \textbf{Value} \\
\midrule
Factorization rank            & $k$                       & $32$ \\
Mean-hyperprior mean          & $\mu_0$                   & $0$ \\
Mean-hyperprior precision     & $\gamma_0$                & $1$ \\
Weight/factor precision prior & $\alpha_\lambda,\beta_\lambda$ & $1,\ 1$ \\
Noise-rate hyperprior         & $c_0,\ d_0$               & $1,\ 1$ \\
Noise-prior shape (empirical Bayes) & $a_0$               & $\approx 1.75$ (fixed after EB) \\
Noise-prior rate              & $b_0$                     & initialised at $1$, learned \\
Main run                      & $n_{\mathrm{iter}},\ n_{\mathrm{burn}}$ & $2000,\ 100$ \\
Retained samples              & $S$                       & $1900$ \\
EB pre-pass                   & $n_{\mathrm{iter}},\ n_{\mathrm{burn}}$ & $100,\ 50$ \\
EB data-rich threshold        & $N_{\min}$                & $50$ \\
Cross-validation              & folds                     & $5$ (grouped by triplet) \\
\bottomrule
\end{tabular}
\end{table}

% =====================================================================
\subsection{Posterior-predictive quantities}
\label{sec:predictive}

For a query $\xx^\ast$ belonging to chemical $c^\ast$, each retained sample gives
a deterministic FM prediction $\yhat^{(s)}(\xx^\ast)$ via Eq.~\eqref{eq:fm}. The
point prediction is the sample mean
\begin{equation}
  \yhat^\ast=\frac{1}{S}\sum_{s=1}^{S}\yhat^{(s)}(\xx^\ast).
\end{equation}
The predictive variance decomposes by the law of total variance,
$\Var(y^\ast\mid\mathcal{D})=\Var_\theta\!\big(\E[y^\ast\mid\theta]\big)
+\E_\theta\!\big[\Var(y^\ast\mid\theta)\big]$, whose two Monte-Carlo estimates are
the epistemic and aleatoric components:
\begin{align}
  \Var_{\mathrm{epi}}(\xx^\ast) &= \frac{1}{S}\sum_{s=1}^{S}\big(\yhat^{(s)}(\xx^\ast)-\yhat^\ast\big)^2,
  \label{eq:epi}\\
  \Var_{\mathrm{ale}}(c^\ast) &= \frac{1}{S}\sum_{s=1}^{S}\frac{1}{\alpha_{c^\ast}^{(s)}},
  \label{eq:ale}\\
  \Var_{\mathrm{tot}}(\xx^\ast,c^\ast) &= \Var_{\mathrm{epi}}(\xx^\ast)+\Var_{\mathrm{ale}}(c^\ast).
  \label{eq:tot}
\end{align}
The epistemic term is the spread of the mean across plausible fits and shrinks as
data constrain the parameters; the aleatoric term is the posterior-mean noise
variance of the chemical and does not vanish with more data. The aleatoric term
depends only on the chemical, so it is constant across species and durations for
a given chemical.

\paragraph{Calibration.} To assess whether $\Var_{\mathrm{tot}}$ is well
calibrated we form posterior-predictive draws that combine both components. For
each held-out observation and each sample $s$ we draw
\begin{equation}
  \tilde{y}^{(s)} \sim \Norm\!\big(\yhat^{(s)}(\xx^\ast),\ 1/\alpha_{c^\ast}^{(s)}\big),
\end{equation}
giving an $S$-point predictive sample. For a nominal central level $\rho$ we take
the empirical $[(1-\rho)/2,\ (1+\rho)/2]$ quantiles of these draws as the
credible interval, and report empirical coverage as the fraction of held-out
observations whose true value lies inside. Coverage is evaluated on the
out-of-fold observations (Section~\ref{sec:cv}), so predictions are scored by
models that did not train on them.

% =====================================================================
\subsection{SSD ensembles and hazardous concentrations}
\label{sec:ssd}

The Bayesian structure lets uncertainty propagate coherently into Species
Sensitivity Distributions. For a target chemical $c^\ast$ at a fixed duration,
and each retained sample $s$, we evaluate Eq.~\eqref{eq:fm} at $\theta^{(s)}$ for
all species $i=1,\dots,M$ and sort the resulting predictions ascending to obtain
one SSD curve $\big(\yhat^{(s)}_{(1)}\le\cdots\le\yhat^{(s)}_{(M)}\big)$ plotted
against the plotting positions $m/(M+1)$. Because every curve comes from a single
coherent parameter draw, the correlated epistemic structure across species is
preserved: the whole curve shifts and rescales as the shared parameters vary,
rather than species scattering independently. The ensemble of $S$ curves is
summarised pointwise by its median and central $95\%$ band.

A hazardous concentration $\mathrm{HC}x$ (the concentration protective of
$100-x\%$ of species) is read off each sampled curve at rank
$\lfloor x\,M/100\rfloor$, giving a posterior sample of $\mathrm{HC}x$ values;
its median and $2.5/97.5$ percentiles are the reported estimate and $95\%$
credible interval. Optionally, aleatoric noise can be injected per species by
adding a draw $\Norm\!\big(0,1/\alpha_{c^\ast}^{(s)}\big)$ to each prediction
before sorting; the SSD figures and HC intervals in the main text use the
epistemic ensemble (no added aleatoric noise) unless stated otherwise.

% =====================================================================
\subsection{Cross-validation and out-of-fold uncertainty}
\label{sec:cv}

Predictive accuracy and calibration are evaluated with grouped $5$-fold
cross-validation. Folds are formed by \emph{grouped} splitting on the
(chemical, species, duration) triplet identifier, so that all replicate
measurements of a held-out triplet are removed from training together and no
replicate of a scored triplet is ever seen during training. Each fold refits the
full model (including the hierarchical noise model, with the EB shape $a_0$ fixed
and $b_0$ learned) on its training partition and predicts the held-out partition;
concatenating the held-out predictions over folds yields the out-of-fold arrays
used for the RMSE, the uncertainty-versus-data analyses
(Eqs.~\ref{eq:epi}--\ref{eq:ale}), and the calibration curve. The full-data model
used to generate the SSDs and the all-chemical HC comparison is fit once on all
$n$ observations with the same settings.

Two practical points. First, the epistemic estimate of Eq.~\eqref{eq:epi} is only
meaningful for chemicals that appear in the training partition of a fold;
chemicals with fewer total observations than the number of folds can have folds
in which none of their rows are present, leaving their chemical-specific
parameters at the prior and deflating their apparent epistemic variance. Such
cold-start chemicals are excluded from the per-chemical
uncertainty-versus-data analyses in the main text. Second, the aleatoric
out-of-fold value for an observation is the sample mean of
$1/\alpha_{c}^{(s)}$ over that fold's retained samples, so it reflects the noise
precision learned from the fold's training data only.

% =====================================================================
\subsection{Computational complexity}
\label{sec:complexity}

Every conditional in Section~\ref{sec:sampler} is a single Gaussian or Gamma
draw, and the residual and factor caches make each parameter update a rank-one
correction touching only the non-zeros of the relevant feature column. One full
sweep therefore costs $O(k\,N_{\mathrm{nz}})$, linear in the factorization rank
and in the number of non-zero design-matrix entries, with an additional
$O(C)$ for the noise block (Eqs.~\ref{eq:post-alpha}--\ref{eq:post-b0}). This
matches the complexity of the unblocked BFM sampler of
\citet{rendle2011bayesian}; the per-chemical noise model adds only the vectorised
$\alpha_c$ draw and one scalar $b_0$ draw per sweep, so the hierarchical
extension does not change the asymptotic cost.




\begin{acknowledgement}
\end{acknowledgement}


\bibliography{references}

\end{document}